% ============================================================
%  R. Nielsen, "Om det oprindelige Forhold mellem Religion og Videnskab"
%  Kjøbenhavn: J. H. Schultz, 1881
%  (Indbydelsesskrift til Kjøbenhavns Universitets Fest, 8. April 1881)
%  TRANSLATION (English) of transcription.tex, printed pp. 1-55
%
%  Scholarly, moderately literal, readable English; American spelling.
%  Nielsen's long periods are broken only where English will not carry
%  them, never across a page marker or a footnote anchor.
%
%  ---- TERMINOLOGY ----
%  Tro / Viden = faith / knowledge; Videnskab = science; Erkjendelse =
%  cognition (Troeserkjendelse / Videnserkjendelse = cognition of faith /
%  of knowledge); Erkjendelsesprincip = principle of cognition;
%  Troesobject / Vidensobject = object of faith / of knowledge;
%  ueensartet = heterogeneous; Objectivering = objectification; Phantasi
%  = imagination; Fornuft / Forstand = reason / understanding;
%  Christendom / Christenheden = Christianity / Christendom;
%  Forestilling / Begreb = representation / concept; opbyggelig =
%  upbuilding.
%
%  ---- CONVENTIONS ----
%  - Page markers \opage{N} and "% --- p. N ---" at the same joints as in
%    transcription.tex; a word the Danish divides across a page turn is
%    kept whole in English.
%  - The same paragraphs, \emph spans (every \emph is letterspacing in
%    the print), footnotes and centred section numerals I.-IV. as the
%    Danish.
%  - Quotation marks „...“ -> ``...'', inner quotations `...'.  Where the
%    print leaves a quotation unclosed or closes one it never opened, the
%    English closes every quotation and records the site in % TR:.
%  - German, Latin and French stay as printed, set in \textit (the print
%    uses roman throughout); Greek as printed.  Bibliographic references
%    in the notes stay as printed ("S." -> "p.", "Anfr. Skr." -> "op. cit.").
%  - Bible references keep Nielsen's Roman chapter numerals, with the
%    English book name.
%  - Printer's defects: the Danish carries them "% sic"; the English
%    renders the evident sense and logs the site in % TR:.
%
%  ---- HOW IT WAS MADE AND CHECKED (2026-09-27) ----
%  Five translators (pp. 1-10, 11-21, 22-32, 33-43, 44-55), each
%  re-reading its English against the Danish sentence by sentence; the
%  pp. 1-8 of an earlier version were revised against the collated
%  Danish.  Mechanical audit (tr_audit.py): per printed page the same
%  \opage joints, \emph spans, footnotes, centre blocks and paragraph
%  breaks as transcription.tex, and an English/Danish word ratio within
%  0.95-1.60.  The only flags are the \textit column, by design (see
%  above).  No independent second reading has yet been made.
% ============================================================
\documentclass[12pt,a4paper]{article}
\usepackage[utf8]{inputenc}
\usepackage[T1]{fontenc}
\usepackage{libertinus}
\usepackage{libertinust1math}
\usepackage[english]{babel}
\usepackage{textalpha}
\usepackage[protrusion=true,expansion=true]{microtype}
\usepackage{setspace}
\usepackage[margin=1.2in]{geometry}
\usepackage{fancyhdr}
\usepackage{hyperref}

\hypersetup{
  pdftitle={On the Original Relationship between Religion and Science},
  pdfauthor={R. Nielsen},
  hidelinks
}

\pagestyle{fancy}
\fancyhf{}
\rhead{Nielsen, \emph{On the Original Relationship} (1881)}
\cfoot{\thepage}

\setstretch{1.3}

\title{\textbf{On the Original Relationship\\[4pt]
       between Religion and Science}}
\author{R.\ Nielsen\\[4pt]
  \small Translated from the Danish by Hans Halvorson}
\date{Copenhagen: J.\ H.\ Schultz, 1881}

% ---- original-page markers ----
\usepackage{marginnote}
\newcommand{\opage}[1]{\marginnote{\footnotesize\textit{[#1]}}}

\begin{document}
\maketitle
\thispagestyle{fancy}

\bigskip

\begin{center}I.\end{center}

\noindent The bitter conflict between religion and science cannot be
resolved unless it is done on the inescapable condition that the
original independence of each is mutually acknowledged. Religion by
itself must then be able to subsist without any science, and science by
itself without any religion; but how is that possible? That religion is
not of scientific origin, ``revealed religion'' no more than the myths,
must surely be conceded; but history teaches that science has scarcely
opened its eyes before it begins to shake the given religion and to
open the conflict. Among the Greeks it was above all the philosophers
who led the attack. Xenophanes pointed out contradictions in the
polytheistic popular faith; against the many gods he set unity, against
their temporal origin eternity, against the inconstancy of their nature
immutability, against their likeness to human beings sublimity,
etc.\footnote{Zeller, Die Philosophie der Griechen, 1ster Theil, 2te
Auflage, Tübingen 1856. p.\ 381.} % *)
In the great poets one traces a gradual change in the nature of the
gods; in Euripides they already begin to fall prey to reflection. A
natural interpretation draws nearer to the cultivated mind. Art held
fast to form, but precisely the perfection of plastic form ate away at
the reality of the gods. The Zeus at Olympia no doubt stood as a
mediator between the popular faith and the philosophical-critical
conception; but the impression was aesthetic, and the overwhelming
aesthetic impression stifled the religious. One \opage{2} may assert
without exaggeration that the scientific thinking which undermined the
authority of the Greek gods had a mighty ally in plastic art.

As regards the relation of the revealed religion of the Old and New
Testaments, the matter stands differently. Endowed with rich
capacities, deep feeling, magnificent imagination, sharp understanding
and strong passions, the people of Israel were nevertheless not, like
the Greeks, especially disposed either to plastic art or to profound
scientific research; in world history they were Jehovah's people, in a
distinctive sense the people of religion, the chosen people of
revelation. An artistic imitation of the divine was forbidden, poetry
stood exclusively in the service of religion, and critically testing
science had no root in the nature of the people. That the Mosaic law,
with its higher ethics, does not permit any thorough scientific
investigation is evident. However much the Ten Commandments may seem to
appeal to human reason, it is nevertheless immediately striking that it
is not the grounds of reason but Jehovah's word on which the emphasis
falls. Let it be well noted that the same God who reveals his will in
the general law also speaks directly in the particular laws.
Concerning the tabernacle Jehovah says (Exod.\ XXV): ``And this is the
offering which you shall take of them: gold and silver and copper; dark
blue and purple and crimson and fine linen and goats' hair'' etc.
``And the tabernacle (XXVI) you shall make of ten curtains \ldots\ The
length of each curtain shall be eight and twenty cubits and the breadth
four cubits for each curtain, one measure for all the curtains''
\ldots\ Let no one say that these special commandments exceed reason's
power of comprehension; on the contrary, the supernatural, what is
incomprehensible to us, lies in this: that the Almighty, the Creator of
heaven and earth, has spoken so precisely about hangings and loops and
has left nothing whatever here to Moses' common sense. It is this that
must seem utterly inconceivable to reason here. In the opposite way,
reason is humbled not only in the Old but also in the New Testament
when the supernatural breaks through so suddenly, so mightily, that
nature cracks, the bond of the laws bursts and the miracle \opage{3}
occurs. Or does the incident of the swine of the Gadarenes perhaps not
stand just as high above reason as the miracle (2 Kings XIII) of the
dead man who fell down into the grave of the prophet Elisha: ``life
came into him, and he rose up on his feet''? In such features there
appears such a rift between divine revelation and human reason that
every reconciliation between religion and science must seem impossible.
If revelation is to hold, then science must go; if science is to hold,
then revelation must be dissolved: a desperate dilemma.

We will not dwell here on particular features of the controversies
which have now for almost two thousand years been conducted between
faith in revelation and reason, controversies in which now religion has
raised the ecclesiastical banner high and been on the point of
subjugating science, now again victorious science, on behalf of
offended reason, has been about to dissolve religion; our task is to go
back to the original ground of the conflict, in order to investigate
whether there might not lie concealed in the ground itself a
misunderstanding, an unclear commingling and confusion of religious and
scientific cognition. Christ's words (John VIII, 32): ``\emph{You shall
know the truth and the truth shall make you free}'', rise high above
all partisan disputes, and must then benefit science just as fully as
religion.

But, it is objected, the truth that is to make us free must necessarily
agree with itself in each and every thing. If something, seen from one
side, is indeed true, but from another side untrue, then this lies not
in the truth but in its presentation and apprehension. Even if all our
cognition is piecemeal, the principle of cognition corresponding to
truth itself cannot possibly be parceled out. Two mutually conflicting
principles of cognition must surely presuppose two truths opposed in
principle and in conflict with each other, but that is altogether
absurd. If religious cognition stands in open conflict with scientific
cognition, this cannot mean that there really are two truths here and,
for cognition, two absolutely heterogeneous principles; rather, one of
the two \opage{4} alleged principles must necessarily be false, either
the religious or the scientific. If the religious principle has
validity, then it must comprise the whole truth, the scientific truth
too; science can then lay no claim to having any principle of its own.
If, conversely, it is the scientific principle that has validity, then
religion is in itself without principle, and must in each and every
thing be subject to the judgment of science. --- We shall investigate
in what follows how matters stand in this respect.

That truth in the eternal sense must be one and indivisible, everyone
will surely concede; but the essence of truth is in itself two-sided,
neither merely something objective nor merely something subjective, but
a unity of both in cognition. Without a cognizing subject there is no
cognition, and without an object that is cognized there is no cognition
either. Truth itself thus can be neither the subjective by itself nor
the objective by itself; the pervading unity of subjective and
objective must run through the whole of existence, embracing the whole
of reality with all its possibilities, in the least as in the greatest.
This man, with his imperfect cognition, must concede; but in conceding
it he concedes at the same time that in his limitation he cannot
possibly see through the content of absolute truth in its countless
forms and concretions. If man is really to cognize the truth, then the
truth must divide itself; and it divides itself like light. As the
light of truth streams into man's consciousness, it deposits two
centers of cognition, two principles, each of which takes on its
relative stamp of the Absolute, namely those of faith and of knowledge;
the first concerns essentially man's own self, his inner being, and is
personal; the second concerns things and their essential connection,
and is to that extent impersonal.

If faith and knowledge are separated from the corresponding principles,
they will constantly flow over into each other. That faith cannot
possibly do without all knowledge is just as evident as that knowledge
cannot do without all faith. The believer must of course necessarily
know what it is in which he believes and what it means to believe;
likewise the knower must believe in his senses and his understanding.
The skeptic, who \opage{5} can believe in neither the one nor the
other, not even in his own skepticism, since the truth is that there is
no truth, is shut off in equal degree from science and from religion.

There is a deep-going difference between religious and scientific
faith, between religious and scientific knowledge. Religious faith has
a psychic-ethical stamp and is essentially a matter of conscience, a
purely personal concern; scientific faith does not concern man in
particular, but expresses confidence in the basic conditions and the
reality of science. Religious faith sets knowledge of God and the
supernatural far higher than acquaintance with the laws of reality and
with natural things: all corresponding knowledge receives its
significance from this alone. The supernatural cannot be sensed, and
must nevertheless make itself known in an immediate way; it then takes
on a sensuous-supersensuous form in imagination, and through this makes
an impression on the feelings: it reveals itself. Athena revealed
herself to Achilles (Il.\ I, 194--98), Jehovah reveals himself to
Abraham, to Moses, to the prophets; even God's revelation in Christ is
no exception here. In the First Epistle of John I, 1 it is indeed said:
``That which was from the beginning, that which we have heard, that
which we have seen with our eyes, that which we have beheld, and our
hands have touched, namely concerning the Word of life''; it might well
seem as though the supernatural were here directly confirmed by the
senses and attested in a sensuous manner; but if we compare with this
(Matt.\ XVI, 15--17), where Jesus says to Peter, who acknowledges him
as the Son of God, ``blessed are you, Simon, son of Jonas, for flesh
and blood has not revealed it to you, but my Father who is in heaven'',
then we come to full assurance that a true knowledge of the
supernatural must be of supernatural origin. The knowledge that has
religious faith as its foundation cannot, then, dissolve faith; when
faith wanes, the supernatural objects are instantly transformed into
myths and fancies, Jehovah and Christ just as well as Zeus and Athena.
In this respect Anselm, with his \textit{credam, ut intelligam},
% TR: "credam" as printed (Anselm's formula is "credo ut intelligam"); logged % sic in the Danish, p. 5
is undeniably right against Abelard, who wanted to reverse the relation.

\opage{6} It is otherwise in science. One may indeed say that scientific
knowledge also begins with faith; but, for one thing, certainty about
what is present to consciousness is not faith but an immediate
knowledge; for another, everything that is provisionally assumed in
faith, because it is assumed in untested confidence, is destined to be
subjected to a testing, whereby faith, according to the Cartesian
\textit{de omnibus dubitandum est}, dissolves into doubt, and then,
when the doubt has been removed by grounding, is transformed into
knowledge. We begin by orienting ourselves in the world; senses,
understanding and reason work together in an immediate way; now it is
the senses, now the understanding, that has a one-sided preponderance,
while reason nevertheless always demands that the reciprocal
determination of what is sensed and what is thought be maintained, for
on this scientific knowledge rests. When investigation begins
concerning the difference between what is changeable in existence and
what, as something unchangeable, must form the foundation of all
changes, science takes its beginning. The task is to investigate the
relation between things and laws. Things are sensed, laws are thought;
but since things can no more exist without laws than laws can be
thought without things, sensation and thinking must originally
presuppose each other in the unity of reason.

That religious and scientific cognition must, according to the
characterization given, be essentially different is then surely
evident. Religious faith, which has its knowledge in imagination and
self-feeling, casts a glance at the natural world of things, assures
itself of its perishability and grasps the supernatural; weighed down
by the law, it turns to the miracles and sees in them testimony that
the free, almighty God is Lord of the laws. The religious laws are laws
of freedom, laws of upbringing, which can be changed according to time
and circumstances, such as the laws concerning clean and unclean
animals; the scientific laws, mathematical, mechanical,
physical-chemical etc., on the other hand belong so essentially to the
nature of things that they cannot possibly be changed without nature
itself having to be changed. All laws of knowledge are laws of
necessity; if they could be broken or annulled, then \opage{7} the
whole world would have to collapse; science must therefore resolutely
attack the miraculous in the name of reason.

There are those, no doubt, who think this too hard a saying. In actual
existence there are certainly nodal points in which the supernatural
and the natural, miracle and law, seem to clash with each other; but
this is only apparent: the new real element introduced by the miracle
must surely, by entering into connection with the natural, direct its
development according to the law. Certainly, we reply, one can
reconcile what is mutually contradictory in the hovering form of
general propositions, but what good is that when the contradiction
shows itself in the fact itself? It is related (Exod.\ III, 2--4) of
the calling of Moses: ``The angel of Jehovah appeared to him in a flame
of fire out of the midst of a thornbush, and he looked, and behold, the
thornbush burned in the fire, and the thornbush was not consumed. Then
Moses said: I will turn aside and see this great sight, why the
thornbush is not burnt up. But Jehovah saw that he turned aside to see,
and God called to him out of the midst of the thornbush and said:
Moses! Moses! and he said: behold, here am I''. How it may be possible
for faith to apprehend this description we can well enough understand,
for faith sees everything in the light of imagination, and imagination
knows how the supernatural combines with the rational, the validity of
the laws with the interruption of the laws; but in knowledge and
science it cannot be apprehended. Human knowledge is in everything,
both the sensuous and the thinkable, so strictly bound to the natural,
the merely rational, that it can neither hear God speak like a human
being nor see how a thornbush can go on burning without being consumed
by the fire. From this, then, it is clear that an object of faith must
be of an entirely different nature from an object of knowledge. If we
place two natural scientists, one on each side of Moses, they would be
able to see the thornbush, as surely as it is an external, real thing,
but they could not possibly either see the angel or hear Jehovah's
word. The thornbush is an object of knowledge, but the angel and
Jehovah are objects of faith.
What can now be inferred from this? That the cognition of knowledge,
the cognition of the natural, must be \opage{8} common to all the
naturally gifted who make rational use of
% TR: the print has "Brug a" (the "f" of "af" dropped); rendered "use of", p. 8
sense and understanding, while the cognition of faith,
the cognition of the supernatural, if indeed such a cognition really
exists, is possible only for ``the elect'', for those who with a
special sense for the supernatural can see what others cannot see, and
hear what others cannot hear.

% --- p. 8 ---
We see that there cannot possibly be a cognition without there being
someone who cognizes and something that is cognized, although the
relation is quite a different one in religion from what it is in
science. That every cognition of the objective must presuppose an
objectification, and every objectification an objectifying
subjectivity, no one will deny who has thought seriously about the
matter. There may be countless objects that no human being has yet
objectified, and there may perhaps be objects for whose objectification
the necessary conditions are lacking; but objects so objective that
they can neither be experienced nor represented nor thought by any
subjectivity whatsoever there cannot be. But if it stands unshakably
firm that no object can be without a corresponding objectification,
then it follows incontrovertibly that an objectifying subjectivity must
extend just as far as the world of objects. Conversely, let
subjectivity be as infinitely clear, as eternal and deep as one
pleases, it must nevertheless vanish like an empty semblance when the
last remnant of the objective disappears. Subjectivity is just as
little able to objectify itself without presupposing an other that is
not subjectivity as so-called pure space or pure time will have any
subsistence when the actual, changeable content is thought away. In
principle, subjective and objective are opposed and yet mutually
belonging determinations, which stand and fall with each other. To
speak of what is objective for man is, in the language of religion as
well as of science, to speak of what can be humanly objectified. The
difference of objects is essentially conditioned by a different
objectification, and the difference of objectification is grounded in
the objectifying subjectivity.

% --- p. 9 ---
\opage{9}Revelation, the objective of religion, can be objectified
only in imagination. This is by no means to assert that what is
revealed is only an imagination. To maintain such an assertion would
in itself be just as absurd as to assert that an object such as a
stone, a rock, was only a sensation, because such things must
necessarily be objectified through sensation. We do not here regard
imagination as a creative faculty, but as a medium of cognition. It is
necessary in science as well as in religion, but its use is just as
different as scientific objectification is different from religious.
In science the interest falls essentially on the objective, in religion
on the other hand first and last on the subjective. A scientific
objectification, e.g.\ of the solar system, presupposes imagination,
but imagination is nevertheless not permitted to introduce a single
member, not even a single feature, unless its reality has been
sufficiently tested in experience and thinking. Although all knowledge,
as we have seen, must necessarily presuppose a subjectivity determined
for the objective,
% TR: the print has a stray comma after "bestemt" ("bestemt, Subjectivitet"); rendered without it, p. 9
it is nevertheless an unavoidable demand in science that the objects
not be colored by subjective light, but be apprehended, without any
admixture of arbitrary elements or irrelevant additions, as they must
be assumed to be in themselves according to universally valid laws.
Whether such objects are pleasant or unpleasant to man, beneficial or
harmful, has no influence on scientific objectification; the objects of
science are in this sense impersonal. It is otherwise with the
religious objects. In religion the objectifying subjectivity must take
account of its own inner being,
% TR: the print has "Frygt." (full stop for comma); rendered as a comma, p. 9
its fear, its hope, its need and passion, its peace and salvation. The
objectification then sets the mind in motion, and imagination imparts
a personal tinge to the objects. In the Psalms of David this feature is
prominent in many ways. In Ps.\ XVIII.\ 7--18 it is said: ``In my
distress I called upon Jehovah \ldots\ Then the earth was shaken, so
that it trembled, and the foundations of the mountains quaked \ldots\
He rode upon a cherub and flew, he soared on the wings of the storm. He
pitched a tent of darkness as his shelter
% --- p. 10 ---
\opage{10}around him. It was his hut, dark waters in thick clouds''.
Here objectification takes place in imagination. Cf.\ Ps.\ LXV,
Ps.\ CIV, Ps.\ CXIV etc.

That religion, even if in another and higher sense, needs the free use
of imagination for its objectification of reality just as surely as
poetry does, can be seen, among other things, from the story of
creation. When it is said in Gen.\ ch.\ 1 that the Spirit of God
hovered over the waters, and that God said let there be light, the
spirit becomes visible in the medium of imagination, and God himself
appears as an object of faith in that he brings forth the light. By
light and its opposite, darkness, is presumably meant sensuous light and
sensuous darkness --- ``it became evening and it became morning'' ---
but the sensuous dissolves into the supersensuous, for God is no
sensuous thing, and the word: let there be light! can be heard only
with the ear of faith. The scientific theory of light cannot discover
the light that is created on the first day and precedes the luminous
bodies, the great and the small light, which come into being only on
the fourth day. Nor can science come to terms with herbs sprouting
from the earth a whole day before the coming into being of the sun, the
moon and the stars, any more than it can catch sight of the extended
firmament that divides the waters above from those below. If faith is
right, then science must be wrong. An endeavor to reconcile faith and
knowledge by letting the six days designate six periods, each of
several million years, is equally ruinous to religion and to science.
By going along with this, religion of course departs from the clear
words of the Bible about evening and morning, and finds no way of
understanding ch.\ 2, v.\ 3, where God hallows the seventh day because
he rested on it from all his work. Nor can science be served by such an
agreement; it lacks an eye for the supernatural, and will not possibly
be able to come to terms with herbs and trees having appeared millions
of years before the sun came into being.

That the religious objectification of things and their connection is
fundamentally different from the scientific, everyone who will take the
trouble to understand the relation must necessarily concede. If we
consider Paradise, the creation of woman and the Fall, Gen.\ II and III,
% --- p. 11 ---
% TR: p. 11 opens mid-sentence; the Danish period begins on p. 10: "Betragte vi Paradiset, Kvindens Skabelse og Syndefaldet, 1. Moseb. II. og III., | vil Forskjellen mellem Tro og Viden vise sig ..." (If we consider Paradise, the creation of woman and the Fall, Gen. II and III, | the difference ... will show itself ...).
\opage{11}the difference between faith and knowledge will show itself, if possible,
still more conspicuously. For the objectification of faith, Eden,
the garden that God planted in the east and in which he put man,
has not merely as much reality as any other garden,
but infinitely more; for here the natural fuses harmoniously
with the supernatural. The river that divides, the gold,
bdellium and the stone shoham are natural, of course, the tree of life
supernatural, and the tree of knowledge at once both natural
and supernatural. Or should there perhaps be nothing natural at all
in the tree of knowledge, how then is it possible that Eve, who
after all cannot be altogether supernatural, can take delight in seeing
it and eating of its fruit? But if the tree were not also supernatural,
how then should it, besides the fruit of knowledge,
be able to bear the talking serpent? That he who, John VIII, 44, is called
the murderer from the beginning must have been in the serpent is
likewise more supernatural than natural. It has, to be sure, not seldom
been declared that the whole account of Paradise and the Fall need not
necessarily be understood literally, but could very well be
taken either allegorically or mythically, without religious faith
thereby suffering any detriment. That this, however, is a concession
dangerous to revelation will readily strike the eye. It
cannot be impressed strongly enough that all faith in revelation is in its principle,
and must be, faith in miracles. Should the miraculous in Paradise really
entitle us to take the whole account as a
myth, then the same would of course have to hold of everything
miraculous in the New as well as in the Old Testament. Christ's
birth of the Virgin, his miracles, his resurrection and ascension
must then dissolve into sheer myths. But when the miracle, in
whatever shape it appears, is myth and revelation
mythology, then religion is subject to the judgment of science, and
% TR: p. 11 the print has a comma after "opløst" before the capital of "Staaer" (so printed, checked at 600 ppi); the English ends the sentence there.
its positive form is thereby essentially dissolved. If, then, one stands
firmly on the principle that our cognition of truth can
have only one principle, since the eternal, absolute truth is only one
and must in all things be equal to itself: then the struggle between religion
% --- p. 12 ---
\opage{12}and science, between faith and knowledge, is a struggle of life and death,
as urgent and decisive as possible.

If it is really true, as faith in revelation teaches, that
the history of mankind proceeds from a paradise, then it can
not possibly be true, whatever natural history asserts, that man
has come into being by developing out of some animal form or other;
if it is really true that Christ is risen from the grave and bodily
ascended --- not to another globe in space, but to
heaven, then there must somewhere be a real and yet super-real,
natural and yet supernatural heaven, in which he dwells, a
heaven that the astronomers do not know. Although all our cognition
is piecemeal, there can nonetheless be no question of parceling out
the truth itself, so that some rare, especially remarkable pieces,
as e.g.\ where the law of gravity is suspended and one walks
on the sea, or where the chemical laws are burst, and water is changed
by a word into wine, belong to religion, while the more ordinary
and natural pieces fall to science. The separation
does not here concern the individual, now surprising, now more
ordinary appearances of the truth; no, what is at issue here is a
principle, the principle of the human cognition of the truth.
If there is only one principle here, then it must necessarily be placed either in
faith or in knowledge. If it is placed in faith, then science must
unconditionally bow to religion; science does not know God
and, however far back it looks, still cannot reach the beginning of
creation. If we place the principle in knowledge, then religion
must subordinate itself to science. One can then say to
him who is not able to appropriate the great results of scientific
research: In your simplicity you may well believe
in Paradise and the Fall, in Christ's supernatural conception
and ascension, in what is awakening and consoling for the common
people in these images and notions; but never imagine that
it is really truth. Truth is only in science, and
he who does not know science never comes to know what
truth essentially is.

That there is only one principle of cognition will, in a practical sense,
% --- p. 13 ---
\opage{13}have the consequence that there can be only one supreme authority
of truth, which must then be sought either in religion or in
science. What follows from this, history has shown. In Catholicism
throughout the whole Middle Ages, indeed even in the old Protestantism
until toward the end of the seventeenth century, the principle
was religious and the overreaching authority on the side of the Church; a
scientific development of form was indeed also in motion, but as
regards the content of truth, reason, darkened by the Fall,
had to borrow its light from revelation, honor it as a
maidservant (\textit{ancilla}) honors her mistress, and resign itself to its captivity
under ``the obedience of faith.'' If the principle of faith is to count as the
only true principle of cognition, then it is perfectly consistent
that the science in which the form of reason is compelled to
develop the content of revelation, namely theology, must be regarded
as the highest and properly true science. As regards other scientific
endeavors, theology is therefore fully within its rights
when, in cases of dispute, it summons them before its tribunal and
judges them by its own laws.

In the eighteenth century the relation begins to turn;
not only in historical-philological criticism, but also in philosophy and
natural science, which, despite strife and disagreement among themselves, were agreed
in recognizing the right of reason, science emancipated itself
from the authority of religion, declared outright that the true principle
of cognition was not religious, but scientific. At first
one sought to purge every trace of the unscientific and to transform
religion itself into the highest science; when the mistake was
discovered, science, by virtue of the solely valid
rational principle, aimed directly at a radical dissolution of all positive
religion. Those who still, in the name of revelation, would make faith
in the supernatural the only true, all-embracing principle of
cognition have dark prospects; for science, with reason
and the natural world as its foundation, is a power, indeed a
great power, that will not a second time let itself be gagged and put in
chains. It would then have to happen by a miracle so great that the whole of
modern culture collapsed into chaos.

% --- p. 14 ---
\opage{14}As regards miracles, the Protestant Church has
hardly been fortunate in its interpretation of the Bible. In the Gospel of
Mark XVI. 15--18, the ascending Savior says: ``Go
into all the world and preach the gospel to the whole creation. He who
believes and is baptized shall be saved, but he who
becomes unbelieving shall be damned. But these signs shall
accompany those who believe: in my name they shall cast out devils, they
shall speak with new tongues; they shall take up serpents, and if
they drink any poison, it shall not harm them; they shall
lay their hands on the sick and they shall be healed.'' No one
will be able to deny that the general Church, from its standpoint, is
right when, in unconditional reverence for Christ's words, it lays decisive
weight on faith and baptism as indispensable conditions
for salvation and blessedness; but it must astonish us in the highest degree
that the Protestant churches, even Bible-firm, orthodox Lutheranism,
take no account whatever of the outward marks with which
Christ himself, in the moment of parting, solemnly designates the believers.
It is a polemic against Catholicism that has misled
the Protestant Church into misjudging the indissoluble connection
between biblical Christianity and its miracles. One
is offended by the Catholic legends of the saints, the many petty, tasteless
miracle tales, and so resorts to the, as it seems, very
reasonable expedient of letting the age of miracles end with the death of the apostles
and pure Christianity continue by the mysterious effects of the Word and the Sacrament
alone, without visible miracles.
In Holberg's Kirke-Historie, Part 1\footnote{Edited by Liebenberg. 1867.} % *)
p. 31, it is said of Christianity: ``No doctrine has been confirmed with so many
and great miracles, although no doctrine needed signs less,
and that in respect of its excellence; for all that Jesus taught
on earth aims solely at the happiness of the human race,
and one sees that its morality has had such an effect that barbarous
lands, even at the times when the doctrine had been obscured by human
inventions, have thereby received another
% --- p. 15 ---
\opage{15}shape. Yet God has, as it were in superabundance, authorized it
not only by the prophecies of the prophets, but by countless miracles,
so that no one's unbelief could be excused.'' --- But who does not now
see that such a separation, which Holberg by no means
invented, but expressly found already in the Lutheran Church, the separation
between a Christianity with miracles and a Christianity without
miracles, is something so altogether unknown in the New Testament that
neither Christ nor the apostles seem to have had even the remotest
inkling of it? By entering upon this separation the
Lutheran Church, despite all protests, without being conscious of the contradiction,
has transferred the principle of cognition from religion over into
science and thereby dug its own grave. It is then incontrovertibly
certain that revelation and reason, faith and knowledge, must wage an
incessant, irreconcilable struggle with each other, until either religion
has subjected all science to itself or, conversely, science has made
itself master over all religion. If one cannot possibly accept
the first supposition without violating science, and cannot possibly
give oneself over to the latter view without violating religion
and attacking the life of the people at its innermost root, then it becomes
necessary to investigate whether the one, eternal, in itself absolute
truth should not, by its reflection in human consciousness,
be compelled to divide itself, to radiate from two different centers
and thereby to deposit two relative, yet for all human cognition
absolutely heterogeneous principles, the one for faith, the other
for knowledge.

That the two spheres of truth, each with its center, cannot
stand in conflict with each other follows from the fact that both have their
origin in the one and same eternal unity of truth; but that
they also cannot fuse together follows from the fact that their
cognition, on account of the separation, must have a different point of
departure, a different goal and a different method. The object of faith
cannot possibly become an object of knowledge, nor conversely. Questions
such as those about creation, Paradise and the Fall, redemption and
restoration cannot be answered scientifically at all, and must
therefore be resolutely turned away in science; whereas all questions
% --- p. 16 ---
\opage{16}about atoms, forces, laws, about the natural formation and
dissolution of things belong to science and do not concern religion
at all. When a teacher of religion sallies forth with a supposedly
decisive polemic against science and would have its criticism of the
biblical doctrine of creation, of the miracles, etc., prove that science
is false, he proves only that scientific
cognition neither is nor will nor can be religious, precisely because
it is and must be scientific. When conversely someone in the name of
science proves that such notions as those of
Paradise and the Fall, resurrection from the dead, etc., must be completely
eradicated because they are all encumbered with intolerable
contradictions: then such proofs have significance exclusively from a
scientific, not from a religious standpoint. Religion can,
with its principle, when it speaks outright, openly and freely declare
that it is and must be unscientific. Neither Abraham nor
Moses nor Christ was a man of science or spoke out of
deeper scientific presuppositions; their grounds are personal,
and the objects of cognition, which are seen in the light of personality,
thereby all without exception receive a personal tinge. The religious person
can in his own way take all human powers into service without
making use of science; there runs through the Old and
New Testament a rich and deep and mighty stream of thoughts, but
they have essentially nothing to do with science. Science
can, according to its principle, declare outright that it is at bottom neither
Christian nor Jewish nor pagan, but universally rational, and
must seek its original independence solely and only in the universally rational.

Only when this fundamental separation, this dualism of temporality
into which the eternal enters, is fully recognized and held fast, can there
be talk of true reciprocal action between the religious and the
scientific in human life. The eternal, absolute truth,
which in its essence is monistic, enters into temporality and becomes
dualistic; man, in whom this dualism comes to clear
consciousness, cannot be rid of the dualism either by ascetic-unscientific
denial of the reality of the temporal or by
% --- p. 17 ---
\opage{17}nihilistic-scientific denial of the reality of the eternal, but must,
in dualistic cognition of both realities, work his way forward in
approximation toward an ever richer understanding of the absolutely true,
concrete monism.

By briefly illuminating the historical misrelations between
religion and science, we shall in what follows attempt to work
our way toward a standpoint from which it can be seen that religion can
just as little be the guardian of science as science the guardian of religion,
but that both, precisely by virtue of the two heterogeneous
principles, can work together in good understanding.

% print: a short centred rule separates Section I from the head; head numeral in heavier type
\begin{center}II.\end{center}

If revealed religion is the truth itself, and the cognition of truth
only one, then reason has no particular principle of cognition.
Revelation must then be infinitely exalted above man's limited reason,
darkened by sin; not even the explanation of the natural order of things
may be left to a reasonable use of the understanding and the five senses alone,
for the senses, after all, extend only to what is lower, and the thinking
left to reason is exposed to great errors. Luther, who otherwise ranks
reason high in all worldly, civil matters, demonstrates its
complete impotence and insufficiency when the talk is of
answering decisive religious questions; in particular, in the
eucharistic controversy he so brushed aside the Carlstadtian reason that he
thereby palpably struck the universally human one. One would then
think that the Lutheran Church, even if by a suitable use of
a subordinate reason it could obtain a well-ordered and coherent
doctrine of religion, could not possibly succeed in obtaining a dogmatics
so scientifically, many-sidedly grounded as the scholastic Catholic one.
And yet it has happened. The Lutherans' confession of faith has to such
a degree developed into science that dogmatics has become the decisive
% --- p. 18 ---
\opage{18}norm for true orthodoxy. Melanchthon, who with his \textit{Loci
Theologici} began religiously and ended scientifically, led the way;
but since he could not properly get the better of reason and in particular
could not open his eyes to the fact that the highest must be precisely what is
contrary to reason, his tendency was seriously combated. Only
in the \textit{Formula Concordiæ} is it expressly laid down how reason
is unconditionally subject to orthodoxy, must humble itself under what is
above reason, indeed bow deeply before what is contrary to reason, and now
the doctrine of religion became orthodox-systematic. Religion again became
science in becoming theology. Like an eagle on mighty wings
the reborn theology rose in the sunlight of revelation to
such a height that it could freely look down on and judge all the
human knowledge that had to keep to the earth. How
matters really stand with this orthodox science we shall
seek to illuminate in a few strokes.

As the all-embracing principle of cognition, faith has Holy Scripture
as its foundation; the first concern, then, is to secure the
infallibility of Scripture by securing its verbal inspiration. It had to be
proved, and was then, if not properly proved, at any rate asserted in faith,
that prophets and apostles have not written a single word from natural impulse
or on their own account, but that it expressly took place according to
the immediate prompting of the Holy Spirit. Every word of the Bible
is then in the proper sense God's word; the instruments chosen for the
writing down are ``God's scribes, God's pens.'' That the authority of
Scripture does not, as the Catholics assert, depend on the authority
of the Church, but on God alone, can be seen from manifold passages of Scripture, John
V, 34; VII, 17; 1 Thess.\ II, 13 and others\footnote{See e.g.\ Brochmand.
Universæ Theologiæ Systema. Articul II. Cap.\ II p.\ 14--15.}. % *)
The purpose of God's having himself willed
to communicate his will to us in writing is solely that we
should learn to worship him in the right way and to conduct our life
so that we can surely attain blessedness\footnote{Op.\ cit.\
Cap.\ I. p.\ 2.}. % **)
The clarity,
completeness and accuracy of Scripture is so great that even the vowels in
the Old Testament must be assumed to be inspired. That
% --- p. 19 ---
\opage{19}no contradictions can possibly be shown between different passages of Scripture
where the same thing is expressly spoken of is self-evident. It
can undeniably go so far that Jehovah is expressly said in one place
to do what Satan does in another, as when it is related in 2
Samuel XXIV, 1 that Jehovah himself tempted David, while
it is said in 1 Chron.\ XXI, 1 that it was Satan who tempted
him; but this can very well be explained if only one is willing to take
reason captive. That Jehovah tempts means, after all, that he cannot
tempt --- ``Let no one say, when he is tempted: I am tempted by God; for
God is not tempted by evil, and he himself tempts no one'' (James
I, 13). Since, then, Jehovah in his wrath wants to tempt David, it is surely
not incomprehensible that he lets Satan do what he himself
cannot; only let one beware of the subtlety of reason that would have
Jehovah himself, with respect to the temptation, be in agreement with Satan.
Nor, as regards the purely historical, can the sacred
authors be assumed, as the Socinians nonetheless assert, to
have erred here and there. To be sure, there are several
difficult passages even in the New Testament. It is indeed a great question,
concerning Luke II, 1--2, whether Quirinius was precisely governor in Syria
when the Emperor Augustus had all the world enrolled for taxation; it is, if
possible, an even more difficult question how the prophet Jeremiah,
Matt.\ XXVII, 9, can be exchanged for Zechariah. If one cannot
take reason captive, one must admit a lapse of
memory; but when faith has become the master of reason, one sees
at once that Jeremiah may very well have said what not he himself,
but Zechariah wrote\footnote{Brochmand. Systema, Part 1,
p.\ 81: \textit{Nobis tutissima responsio est hæc: verba illa qvæ a Matth.\
adscribuntur Jeremiæ, a Jeremiæ qvidem dicta esse, (non enim dicit
% TR: Latin kept as printed, with the print's defects "al Jeremiæ" (for "a Jeremia"), "Zaeharia" (for "Zacharia") and "græccum" (broken "græc-cum" over the line); the Greek word copied verbatim.
Evangelista, hoc al Jeremiæ scriptum, sed}
ῥηθὲν \textit{ceu dictum est): verum a
Zaeharia in litteras sacras relata. Nam aliæ qvæ afferri solent
responsiones, aut Textum græccum corruptelæ arguunt, aut Matthæo
dicam lapsus memoriæ scribunt, aut conjecturis plane incertis
nituntur.}}. % *)

Since the Bible, however, is subject to interpretation, and every
interpreter is driven by his spirit, there arises the serious question:
% (Danish word "Spørgsmaal" divided across pp. 19/20; English word kept whole on p. 19)
% --- p. 20 ---
\opage{20}by which interpretation, then, is one to know the true
Spirit? Luther's declaration: ``\textit{Der heilige Geist ist der allereinfältigste
Schreiber und Redner, der im Himmel und Erden ist,
darum auch seine Worte nicht mehr denn einen einfältigen Sinn
haben können, welchen wir den schriftlichen oder buchstablichen
Zungen-sinn nennen}''\footnote{Walch XVIII 1602.} % *)
gives us a provisional guidance. Only in
the spirit of the Bible can the Bible be interpreted; the spirit of the Bible is the Holy
Spirit, consequently the interpreter must have the same Holy Spirit
if his interpretation is to be true. When Jesus was baptized by
John, it is said: ``then the heavens were opened above him,'' as he
went up out of the water, ``and he saw the Spirit of God descending like a dove
and coming upon him.'' And behold, from the heavens there came a voice,
% TR: p. 20 the print closes the quotation after "Velbehag" but never opens it (logged sic in the Danish); the English supplies the opening mark before "this is".
which said: ``this is my Son, the beloved, in whom I am
well pleased.'' Matt.\ III, 16--17; in Luke III, 22 it is expressly
said that the Holy Spirit descended in bodily shape like a
% TR: Greek copied verbatim, with the print's "πνῦμα" (for πνεῦμα), unaccented "το" and "ἔιδει" (for εἴδει).
dove (καταβῆναι τὸ πνῦμα το ἄγιον σωματικῷ ἔιδει). If one is
now to take this quite simply according to the letter, then a clear and
definite notion of how the Spirit looked in a
real dove-shape is difficult, for he who looks simply at
the dove cannot see the Spirit, and he who fixes his gaze on the
Spirit cannot see the dove; it becomes equally difficult to understand
whether it was Jesus alone who saw the dove and heard the voice, or
John as well. When one reads the Gospel of John I, 15--32,
III, 27--36, one would indeed think that the Baptist himself had both
seen the dove and heard the voice; but if we compare with this Matt.\ %
XI, 2 and 3, and Luke VII, 18--35, one is tempted to doubt
that the Baptist was really in agreement with himself about what he
had actually seen and heard.

The relation between spirit and letter is by no means so
comprehensible as it might seem at first glance. If we begin with
the literal meaning, then reason and understanding have so much
to do with the meaning of words and with context, with grammar and
hermeneutics, that exegesis, taken scientifically, easily forgets the
% --- p. 21 ---
\opage{21}Holy Spirit; and yet it is the Spirit that gives understanding its
right stamp. That there is spirit in the Catholic Church just as
well as in the Protestant ones, surely no one will deny. These churches
lie in conflict with one another precisely because each of them claims
to have the true spirit of revelation. Catholics, Zwinglians,
Calvinists each claim for themselves to have the Holy Spirit in preference
to the Lutherans; how is this conflict now to be settled by
invoking the Spirit? We hear Luther inveigh against the false prophets:
``\textit{Schrift, Schrift, Schrift aber ist nicht eitel Geist, davon sie geifern,
der Geist müsste es allein thun; die Schrift sei ein todter Buchstabe,
und könne nicht das Leben geben. Rühme aber nicht viel
vom Geiste, wenn du nicht das aeusserliche offenbare Wort hast,
denn es wird gewisslich nicht ein guter Geist sein, sondern der
leidige Teufel aus der Hölle}.'' We can agree with Luther in this,
and yet get not a hair's breadth further; his words to the Zwinglians
at the Colloquy of Marburg 1529: ``\textit{Ihr habt einen andern
Geist denn wir},'' show clearly enough what the invoked unity of
body and bread, of blood and wine, amounts to. If the decisive thing is to lie
in faith, then everyone must of course abide by his faith; the Zwinglian
then has the very same right as the Lutheran. Even if the words
of institution: ``this is my body, this is my blood,'' are taken and used
quite according to the letter, there is nevertheless in the letter itself no
understanding; the understanding, which comes with the interpretation, must necessarily
become either reasonable: ``this signifies,'' or contrary to reason:
``this \emph{is} what it nevertheless cannot yet be, either naturally or
supernaturally.'' When Christ distributed bread and wine, he
still had, after all, his own natural body and blood; if the bread and
the wine are \emph{not merely to signify}, but in a \emph{supernatural manner
to be} the natural body that is to be given up, the natural blood
that is to be shed on the cross, then we must necessarily get either
an altogether incomprehensible ambiguity, or a mystery so
impenetrable that faith itself must cling to the absurd ---
\textit{credo qvia absurdum est}. But the absurd gives not only no
scientific, but also no religious understanding; it halts
every possible apprehension by the absolutely meaningless. To a
% TR: SLICE ENDS mid-sentence (continued at the head of p. 22). Full Danish sentence across pp. 21/22: "Til en | saadan Conseqvens har den lutherske Kirke ikke kunnet hengive sig; den er bleven staaende ved Troen, d. e. ved \emph{sin} Tro paa Mysteriet." This part renders "Til en" as "To a"; the next part should continue: "consequence of this kind the Lutheran Church has not been able to give itself over; it has remained standing by faith, i.e. by \emph{its} faith in the mystery."
% --- p. 22 ---
\opage{22}such a consequence the Lutheran Church has not been able to surrender
itself; it has remained standing by faith, i.e., by \emph{its} faith in the
mystery.

Now if each particular church community must be conscious of having a
particular faith, then it must of course also have particular formulas for
its distinctive confession of faith, particular symbols. When the faith of
the church is then to be built solely and exclusively on Holy Scripture, the
validity of the symbols must necessarily rest on their agreeing precisely
with the doctrine of Scripture. The Lutheran Church does not claim that its
symbols are infallible; it holds to them only \emph{because} (\textit{qvia})
and \emph{insofar as} (\textit{qvatenus}) their content is in essential
agreement with the content of Scripture. But here the question of
interpretation comes up again, for the movement goes in a circle. An
interpretation without a particular confession of faith can be reasonable,
acute, profound, and yet not true, because the Spirit of truth is lacking;
but the Spirit of truth is the Spirit of faith, and true faith lives only in
the true Church. For the Lutheran, the spirit of the Lutheran Church must be
one with the Holy Spirit of revelation; the Lutheran must thus understand
the symbols in the spirit of the Bible, but from this it follows conversely
that he must interpret the Bible in the spirit of the symbols; if he
deviates from this rule, i.e., if he breaks the circle, he becomes a
heretic.

If we now cast a glance at the dogmas, it is not what is peculiar to the
individual churches but what is common to them all, the universal, on which
we shall dwell in this connection. The Fall is brought about by a temptation
of the Devil; the restoration is owed exclusively to the merit of Christ. To
demonstrate the scientific character of dogmatics as a whole, we therefore
need only cast a glance at the relation between Christ and the Devil in
Lutheran dogmatics. Whether Christ would have become incarnate if man had
not sinned (\textit{atiamsi homo non peccasset}), a question which Thomas
Aquinas treats at length\footnote{Gent.\ IV c.\ 53, 55. Cf.\ Dr.\ Karl Werner. Der heilige Thomas von
Aqvino 1859. 2ter Band. p.\ 620.}, % *)
% TR: "atiamsi" (for etiamsi) kept as printed (Latin stays as printed); the sense is "even if"
% --- p. 23 ---
\opage{23}the Lutheran Church has not undertaken to decide\footnote{Brochmd.\ p.\ 725--26.}; % *)
but it has instead investigated in great detail the relation between
Christ's divine and human nature, and has thereby made it quite clear how
He who was revealed in order to ``\emph{destroy the works of the Devil}''
really had the power to complete his mission. Although the Devil's power
was really broken when Christ revealed himself, the Devil's kingdom
nevertheless continues to the last to resist the kingdom of Christ.
Luther's faith in a real, personal Devil was just as unshakable as his
faith in a real, personal Christ. Luther agrees with Tertullian that not
only evil thoughts and desires, but also crop failure, pestilence, fire and
other external misfortunes have the Devil as their author. In the
\textit{Cathechismus major} and in many other places Luther seeks in the
most urgent manner to impress upon us how persistent, many-sided and
untiring the activity of the devils and demons is, so that we may be on our
guard. ``\textit{Die Heiden},'' he says, ``\textit{wissen nicht, woher das
Unglück so plötzlich kommt; aber wir wissen es, das es eitel Teufels Arbeit
ist, der hat solche Helleparten, Bleikugeln und Büchsen, solche Spiesse und
Schwerter, damit er unter uns schiesst, wirft und sticht, wenn Gott es ihm
erlaubt \ldots}''; and in another place: ``\textit{Ein Christ soll wissen,
dass er mitten unter den Teufeln sitze, und dass ihm der Teufel näher sei
den sein Rock oder Hemde, ja näher denn seine eigene Haut, dass er rings um
uns her sei, und wir also stets mit ihm zu Haar liegen und uns mit ihm
schlagen müssen}''\footnote{Cf.\ Walch.\ X, p.\ 1234, XIII, p.\ 2550.}. % **)

It is, however, not Luther's personal authority but that of the Bible,
namely the fundamental doctrine of the New Testament, that makes it
necessary for Lutheran dogmatics to build, if we may say so, on belief in
the Devil, and to follow the arch-enemy's struggle against the kingdom of
Christ until the Day of Judgment, when he himself with all his followers
shall be bound forever in hell. The Lutheran dogmaticians are for that
reason just as careful to investigate how far the Devil can actually
% --- p. 24 ---
\opage{24}go, and on what his power rests, as to develop faith in the merit of
Christ\footnote{See e.g.\ Brochmand ll.\ Tom.\ I p.\ 250. Qvæst.\ XIII. An Diabolus
cogitationes cordis novisse posse. Pg.\ 251. Qvæst.\ XIV. An Diabolus vera
miracula edere qvæat. Pg.\ 253. Qvæst.\ XV. An Diabolus
animas demortuorum in hanc vitam revocare, et substantiales rerum
formas immutare possit? etc.\ etc.}. % *)
In this matter we must grant that Strauss is right when he maintains ``that
the idea of the Messiah and his kingdom is as little possible without a
kingdom of demons, with a personal head as its counterpart, as the north
pole of a magnet without a south pole''; we must concede to him that ``the
pious narrow-mindedness which feared to lose Christ if they lost the Devil
has seen far more correctly than Schleiermacher with his postulate that
belief in the Devil could in no way be set up as a condition for faith in
Christ''\footnote{Die christliche Glaubenslehre von Dr.\ D.\ F.\ Strauss. Zweiter Band.
1841. p.\ 15.}. % **)

The orthodox Lutheran dogmatics, whose flowering falls in the seventeenth
century, must surely be regarded as an attempt, serious of its kind, to
present the undistorted content of the Bible in strictly scientific form, a
truly magnificent attempt; but it was bound reason that, as the slave of
faith, had to carry out the whole work; reason thus gradually became
conscious of its powers, drew the principle of cognition over to itself
little by little, and thereby prepared its complete emancipation. The Devil,
who with his demons had mistreated it above all, was now himself to furnish
the footstool on which it stood when it began to rise to its full height. A
Reformed theologian in Holland, Balthasar Bekker, a zealous adherent of the
Cartesian philosophy, published a highly sensational work, ``The Bewitched
World'',\footnote{De betoverede weereld, Amsterd.\ 1691--94. In German: Die bezauberte Welt.} % ***)
in which he combats belief in the influence of evil spirits, demonic
possessions, witches and the like. Relying on reason, he demonstrates the
ignorance in which the many absurd accounts of the Devil and his
manifestations in paganism, Judaism, Mohammedanism
% --- p. 25 ---
\opage{25}and Christianity had to be assumed to have their origin. The narratives of
the New Testament about the demoniacs he explains in a natural way, and he
assumes that Jesus himself only meant to reform the moral element in faith
without disturbing ingrained notions. On the story of the temptation he
expresses himself, taking the Fall as his starting point, among other
things as follows: ``\textit{Denn vor fest und ungezweiffelt anzunehmen,
dass ein geschaffener Geist, und der von Gott verworfen ist, auff des
Menschen Seele oder Leib wircken könne, davon ist schon hier vornen
angezeiget, wie viel daran lieget. Hat der Satan eusserlich solch Gespräch
mit dem Menschen geführet, da die Welt erst den Anfang hatte, wie lieset man
das nicht mehr? Selbst nach Verlauff von viertausend Jahren in der
berühmten Unterredung mit dem andern Adam, da der Teufel zwar ausdrücklich
genennet wird, ist das nicht geschehen. Denn obschon beydes Lucas und
Mattheus das umständlich erzehlen, so muss man sich dennoch wol hüten, dass
man nicht alles nach dem Buchstaben verstehe. Oder man müste zugleich
glauben, dass Moses und Elias beyde eben so wol auff dem Berge persönlich
gewesen und sich mit dem Seligmacher unterredet haben, wie solches ein
andermahl, ja dreymahl wird erzehlet. Matth.\ XVII, 3. Marc.\ IX, 4. Luc.\
IX, 30, 31. Gleichwol war Moses nicht leiblich wie Elias in den Himmel
auffgenommen worden, sondern gestorben und gar gewiss begraben, denn Gott
hat es selbst gethan. Deut.\ XXXIV, 6. War er nun denn von den Todten
auferstanden und leibhaftig gen Himmel gefahren: Wunder, dass die Schrifft
ein so grosses Wunder nicht meldet. Wer denn von solchen, als im Gesichte
geschehen ist, so redet, als etwas das nach dem Buchstaben verstanden wird,
der weiss, wie Petrus, auch nicht was er saget}''.\footnote{Op.\ cit.\ p.\ 135--136 ff.} % *)
% TR: the note *) on p. 25 has no call mark in the print (logged % sic in the Danish); anchored, as there, after the quotation it cites

But if we see from this that reason cannot tolerate the Devil and therefore
falls upon his flanks wherever he shows himself, then it also becomes
intelligible to us that free reason, once its eye has been opened and
criticism has properly begun, cannot confine itself to pursuing the Devil,
but, taking its bearings over all the Bible's
% --- p. 26 ---
\opage{26}domains, must have the right to consider how it can come about that Moses,
who was buried by God himself, and Elijah, who went up alive to heaven,
could appear on one and the same mountain; it is, when we look more
closely, not the Devil alone but the miracles that are an offense to free
reason. One will likewise see that reason, once it has got hold of the
principle of cognition, cannot be different in the different churches, but
must rise above the confessional limitations. That the Reformed Church has
a different confession of faith from the Lutheran is certain, but who would
infer from this that reason itself must be different for the Lutherans than
for the Reformed? The Reformed pastor was deposed for his ``\textit{Bezauberte
Welt}'', but his book nonetheless found acceptance among all reasonable
people in the Lutheran (Semler and Thomasius) as well as in the Reformed
Church. In Denmark not only has Holberg's wit made a considerable clearance
among Satan's papers, but the rationalism that began with him has, if with
slow and cautious steps, approached the altar.

On the transition from orthodoxy to rationalism stands Christian Bastholm, a
thinker unprejudiced for his time and respectable. Bastholm's reason is
finished with the Devil to such a degree that he is not even expressly
named in the dogmatics\footnote{Den Christl.\ Religions Hoved-Lærdomme. Kiøbhvn.\ 1783.}, % *)
% TR: the call in the print is *) but the note at the foot is marked **) (logged % sic in the Danish); no consequence for the English
p.\ 265, where the story of the temptation is touched on. ``The seducer of
the world,'' it says, ``presents himself to him'' (Christ); ``many and great
trials are put to his innocence; but he wins the victory over them all. This
temptation was necessary in several respects; before he began to associate
with men, he had to come to know man, his frailties and his afflictions, so
that his own experience might instill in him tender feelings for the misery
of men and their weaknesses'' \ldots\ One would, however, be notably
mistaken if one concluded from this that Bastholmian reason was now also
agreed with itself to pass sentence on everything supernatural and to reject
the miracles. As
% --- p. 27 ---
\opage{27}reason sneaks the principle of cognition over to itself, it does so with
such caution, not to say dexterity, that the less attentive reader notices
almost nothing. Of ``Jesus' Miraculous Works'' it says, among other things,
p.\ 307 ff., ``that they exceeded all the forces of nature and suspended its
best-known laws \ldots\ We do not at all fear that anyone will ever carry
his knowledge of nature so far that he shall put the works of Jesus to
shame, that he should be able to heal the sick merely by a command, merely
by the laying on of a hand, or merely by touching him, or that by his mere
will he should heal those who are even absent. We do not at all fear that
anyone should ever bring his insights into nature to such a height that
merely by rebuking the waves and the storm winds he should stop them in the
midst of their raging, or, by calling to the fish of the sea, should gather
them in the same instant. If it were ever possible to do by the forces of
nature the same works that Jesus Christ has done, then one ought surely to
be able to find some traces of this hope already in our times, after the
great strides that have been made, and the height that has been reached, in
all the natural sciences. But of this there is not even the slightest
trace.'' ``But in this there is of course nothing unreasonable; if Christ
really is the one he said he was, the revelation of divine truth, then it is
of course entirely reasonable that he has been able to do works which
otherwise surpass all human powers.'' ``This inference,'' it says, p.\ 318,
``is so simple that the simplest can grasp it. On this our reason grounds
its tranquility and its conviction; more it need not know.'' In the
doctrine of ``The Apostles' Miraculous Works'' it says, p.\ 404, 405 ff.:
``By a wonder we understand a change in nature which either exceeds or is
contrary to the forces of nature and its laws, at least at the time when it
occurs \ldots\ That such wonders are possible, no one can doubt who knows
the omnipotence of the infinite Being. The laws of nature are not in
themselves necessary; they are only emanations of the free will of the
Creator, who has arranged nature as it is, has granted creatures the powers
they possess, and has thus
% --- p. 28 ---
\opage{28}connected causes with their effects as we now find them. But he who has
given a law can also suspend it, when wise purposes demand it'' \ldots\ ``If
the Creator once by free choice gave fire the power to shine and to burn,
then he can also deprive it of this power again, or hinder its effect''
\ldots\ All this now looks quite reasonable. How reasonable the whole of
revelation has become in Bastholm's conception can be seen, among other
things, from the chapter on ``Jesus' Ascension'', p.\ 369--382. ``After Christ
had walked about for 40 days with his disciples \ldots\ he left the earth
with his visible presence, in order to take possession of heaven \ldots\ It
is the order which the Lord's wisdom has established, that after the
tribulations of life follows death, after death the resurrection, and after
the resurrection an eternal possession of heavenly bliss in soul and body.
Now Christ had died, he had risen; nothing was therefore more natural than
that he should now take possession of the blessings of heaven \ldots\ Reason
therefore finds nothing unreasonable in the doctrine that Jesus Christ
\ldots\ was taken up into heaven, in order to take possession of the glory
and bliss that belonged to him as God-man \ldots\ If one can here, as with
the resurrection of Christ, prove that it was impossible either that the
apostles could themselves be deceived in this event, or that they could
wish to deceive us with a false testimony, then the faith that we place in
their report of it is well founded and reasonable.'' Christ, in the author's
opinion, used four means to convince the apostles and posterity that no
deception had taken place, that what was seen here had not been a delusion
of the senses and imagination, but full reality. ``First, this was a piece
of Jesus' wisdom, that he did not vanish among them. After his resurrection
he was sometimes wont to appear among them quite unexpectedly, and again to
vanish before their eyes just as unexpectedly; perhaps he wished thereby to
convince them that the body which he had brought out of the grave at the
resurrection had other and higher properties than the body that had been
laid in the grave. But now, when he would leave the earth, another purpose
demanded another procedure \ldots\ Our Redeemer was slowly and visibly
taken up into heaven \ldots\ The
% --- p. 29 ---
\opage{29}second piece of Jesus' wisdom was this, that he chose the Mount of Olives
from which to make his ascension. No place could be more convenient for this
purpose. This mountain lay a Sabbath day's journey from Jerusalem; they thus
had opportunity enough to see and observe Jesus in the free and open air. He
spoke with them on the way, gave them his commands, repeated his promise of
the Holy Spirit, whom he would send them, and blessed them \ldots\ This
mountain was also the highest mountain around Jerusalem; here there was thus
a free view and a purer air, which afforded them a convenient opportunity to
behold so important an event quite unhindered.'' The third piece of Jesus'
wisdom must strike us in particular as bearing a predominant stamp of
reasonableness. ``One asks why Christ did not here assemble all his enemies,
in order by so great a wonder to convince them all at once of their unbelief
and of his divine mission? I could answer much to this question; I could
answer that they would have let themselves be convinced by this wonder just
as little as by all the others \ldots\ The Jews ascribed his other miraculous
works to the Devil; they could thus say the same about this one. When a man
is hardened to such a degree, proofs are of no more use \ldots\ I therefore
see no reason why Christ should ascend in the sight of his enemies; but I see
an important reason why he should not do so. Had Christ publicly announced
at what time and at what place he would ascend into heaven, he would
undeniably have drawn all the inhabitants of Jerusalem after him; for who
does not feel the power of curiosity, and who would not walk a little way to
be a spectator of so extraordinary an event? But what noise, what tumult and
disorder must there not necessarily have been in so numerous a swarm, where
one pushes the other, where each wants to be the nearest, and one drives off
the other, as on a stormy sea one wave drives the other.'' Surely something
of the most reasonable that can be said in a scientific dogmatics.

In Bastholm the principle of reason still lies, if we may say so, swaddled
in the supranatural and kicking in its swaddling clothes; but in
% --- p. 30 ---
\opage{30}German literature reason had already earlier become conscious of its
superiority. The Wolffian philosophy so thoroughly permeated the Wertheim
Bible translation\footnote{Erster Theil, 1735.} % *)
that the Spirit of God cannot hover over the waters, but ``a violent wind''
must begin to blow; the concepts are defined so sharply that for the
understanding of Lev.\ XVIII, 7 it must even be added, for clarity's sake,
what a mother is, namely ``a woman who begets children in company with a
man''. The bookseller and author Nicolai in Berlin published from 1765 to
1805 the ``\textit{Allgemeine}'' and ``\textit{Neuallgemeine deutsche
Bibliothek}'', in which reason became quite popular. The theologians
themselves worked along at the enlightenment of Germany. The Provost of
Cölln, Wilh.\ Albr.\ Teller, gave a public reply to an open letter from
Jewish heads of household in Berlin, in which he declared himself willing to
acknowledge them as Christians on the basis of their mere deism. Greater
liberality one could hardly demand; it is as if we heard the very voice of
reason. ``In the year 1833, e.g., the Evangelical General Superintendent
Röhr's `\textit{Predigerbibliothek}' called the doctrine of all Christian
churches concerning the triune God an `\emph{antichristian dogma}' (1833,
H.\ 1. p.\ 11), and called Augustine, on account of his doctrine of sin and
grace, and because he had not remained unconverted, one `who desecrated the
Gospel, and a dissolute man who had become one of the saints' (ibid.\ p.\
13), but Pelagius (p.\ 15) `the venerable defender of reason against
unreason, a man who already in his own time had the satisfaction of seeing
the wisest and best stand on his side', (calling the Evangelical Church
fortunate (p.\ 16) if only Luther had not merely reformed back to the father
of his order, Augustine), and declared (p.\ 48): `Supranaturalism leads away
from the worship of God to the worship of idols' etc.; indeed, in that same
year 1833, the above-mentioned high clergyman, and that not merely
\emph{privatim}, but quite publicly and with \emph{ecclesiastical force of
law}, presented the foundation of the whole of Christianity, the word of the
cross, the doctrine of Christ's death, `a death of redemption and atonement
\emph{for the good of men} (``for the expiation of human wickedness'')
% TR: Danish word order ("en til Menneskenes Bedste (...) udholdt Forløsnings- og Forsoningsdød") inverted so that the p. 31 joint and the \emph span stay in place; the letterspaced "en" (a) is left unemphasized
% --- p. 31 ---
\opage{31}endured', as `a consequence of the total corruption of the pure Gospel', as
`\emph{a fiction of ignorant and hair-splitting church teachers}'. Indeed,
one went still further --- from Halle the fall of all dogmas was proclaimed
(\textit{Halliche Allg.\ Lit.\ Zeitung}; März 1833 Ergänz.\ Bl.\ p.\ 234,
decidedly expresses the hope that `\emph{when all Christian dogmas fall,
sooner or later}', all men would at last unite on a Christian
morality).''\footnote{H.\ E.\ F.\ Guerike. Haandbog i Kirkehistorien. Transl.\ by Sverdrup.
Christiania. 1842. Anden Deel. p.\ 338 note.} % *)
% TR: the print closes the outer quotation (opened at "I Aaret 1833") inside the parenthesis, "Moral“)", and sets a full stop both before and after the call, "Moral“).*)." (logged % sic in the Danish); English closes the quotation after the parenthesis and keeps one full stop

Vulgar rationalism was, however, too superficial to be able to satisfy
profound spirits; for that a speculative reason was required, one capable of
seeing through the identity of the supernatural with the natural.
Dogmaticians such as Daub and Marheinecke looked down on the rationalists,
but up to Schelling and Hegel. Here one found a development of reason so
deep and so wide-ranging that it held the whole of Christianity. When
representation had been dissolved into concept, theology all at once became
so orthodox that the Devil himself went along into the concept. When reason,
by making its principle of knowledge into the principle of revelation, was
on the very point of devouring the whole positive content of orthodoxy, it
became so speculatively swollen that it was close to losing its
understanding. Things might really have looked grave for its precious life,
had one not come in good time to the aid of the swollen patient with a
radical and cleansing remedy; this remedy was a purgative criticism, applied
above all by Dr.\ Strauss.
% TR: "stod i Begreb med at" (was on the point of) plays on "Begreb" (concept) in the sentence before; the pun is not reproducible

However much one may regret that Strauss has been blind to the divine
essence of religion, one must nonetheless, on behalf of science, be grateful
to him that his critical treatment both of the life of Jesus and of church
dogmatics, even if there may be a good deal to object to in the details, has
shown that religious faith must perish when it is to be justified in
knowledge, i.e., when religion is to be conceived and grounded
scientifically. ``Strauss's talent consists in uncovering
% --- p. 32 ---
\opage{32}the confusion and annihilating the illusions; but his criticism is only
\emph{dissolving}, the result remains \emph{negative}.'' ``His dogmatics is
no dogmatics, but only a critique of the individual dogmas, a repertory of
dogmatic representations. Despite all neatness in the external ordering and
delimitation of the material, and despite the sureness of the tendency of
understanding, an enormous lack is nevertheless perceptible, and the
nihilistic background produces a feeling of disconsolateness and
emptiness''\footnote{Dr.\ Karl Schwarz. Den nyeste Theologies Historie. Danish
translation. Kjøbenhavn. 1879. p.\ 779.}. % *)
This view must surely be approved by all who still expect a new theology
from science, but it is precisely this expectation that the Straussian
criticism is expressly designed to annihilate. That he has accomplished his
task in this respect is entirely beyond doubt. One may remodel rationalism
as much as one likes; the rational principle will never become religious.
K.\ Schwarz commends Hase's rational theology. ``The principle governing
everything in his dogmatics,'' he says\footnote{Op.\ cit.\ p.\ 413.}, % **)
``was of course man's \emph{relative freedom}, and the love for the kingdom
of God grounded on it. The Trinity that he developed was also quite
different from Schelling's, which turned into a physical world-event, a
theogonic process, whereas with him it was ethical and was led back to the
practical-biblical content. In general he asserted, in full consciousness,
the ethical character of his theological system, and grounded it on man's
moral-religious freedom in his independence before God, and on his love and
infinite striving for communion with God, as on the firmest foundation. This
ethical principle is indeed the imperishable kernel of truth in
rationalism.''

But the ethical principle is of course not a specifically scientific
principle. Reason can postulate the reality of freedom and recommend to men
in practical life that they assume it and believe in it; but it is not with
freedom as it is with natural things. Everything that science is to be able
to prove, it must prove in virtue of either absolute or conditional
necessity; but
% --- p. 33 ---
\opage{33}when freedom itself is grounded on necessity, then scientific
reason can with full right doubt its reality. Still less can love of
God ground any scientific assurance of God's love, of God's
personality, any more than a personal need for an eternal life can
furnish scientific proof of the immortality of the soul. That reason
with its principle should really ground a doctrine of the Trinity
which scientifically led back to the ``practical-biblical content'' is
an assertion that betrays an utterly unclear confusion of the
character of faith and of knowledge, and constantly mistakes religious
representations for scientific facts, thoughts of revelation for
universal thoughts of reason.

The God of religion, however one turns and twists the matter, is, as
the God of freedom and of love, a supernatural being. Now if it can
truly be said that reason, in all its scientific character, has no eye
for the supernatural, then neither has it any eye for the God of
freedom, of personality and of love. Reason must, passing over the
miraculous---indeed, insofar as it is to make its principle valid,
denying the miraculous---in creation as well as in the incarnation and
in inspiration, hold strictly to nature and its laws. Where, then,
will ``rational theology'' go with its love of God and its exalted
ideal of freedom? It must seek its deity in nature and bind it to the
laws of nature; it must give up freedom and deify necessity. --- He
who would pay homage to Comte's religion is, after all, compelled
either to go beyond reason at a venture, or to turn back and bow to
Littré's criticism.

\begin{center}III.\end{center}

That religion cannot appropriate the principle of cognition in its
whole extent and thereby draw science under itself has been
% --- p. 34 ---
% TR: "uu" (turned n, logged sic in the Danish) read as "nu" = now
\opage{34}shown in the foregoing; the question now becomes whether
science with its principle of reason can then master religion, or
whether, perhaps, by pursuing its own goal and discovering an
essential limit, it will come to see the necessity of granting
religion a distinctive principle of faith. To see how matters stand
in this respect, let us cast a glance at the paths of science. It has
fought, and with its principle must still constantly fight, its way
forward through a double series of oppositions, the one indispensable,
the other in a high degree obstructive and troublesome. Among the
indispensable oppositions we count all those that spring from a
genuine interest in the advancement of science, objections by whose
answering errors committed are removed and method is secured; among
the hampering, troublesome and encumbering ones, on the other hand, we
must count the oppositions and objections that come from the
uninformed, who from one interest or another would gladly bring science
to a halt. It is, for the sake of the principle, especially to the
latter that we must direct attention.

Nicolaus Copernicus, who already in 1530 had his new world
system\footnote{De Revolutionibus orbium coelestium, Norimbergæ MDXLIII.} % *)
finished, nevertheless let it be printed only in 1543, and touched the
first copy with a dying hand. That it was not so much fear of the
polemic of the Ptolemaic astronomers as, far more, of that of the
theologians and priests that held him back can be seen from his
dedication to Pope Paul III, in which he foresees that empty babblers
(ματαιόλογοι), devoid of all mathematical knowledge, will presume to
pass judgment on the work by willfully distorting some passage or
other in Holy Scripture (\textit{propter aliqvem locum scripturæ male
ad suum propositum detortum}); and this presentiment was indeed
fulfilled abundantly enough. Tycho Brahe's observations, Kepler's
laws, Galileo's discovery of the moons of Jupiter, and finally
% TR: "1689" as printed (the Principia appeared in 1687); kept, since the error may be the author's
Newton's mathematical Principles (1689) mark magnificent advances in a
period in which church orthodoxy otherwise held the sciences in
chains. No doubt the theologians must have been embittered at these
ungodly theories, for although they
% --- p. 35 ---
\opage{35}did not properly understand them, this much was clear, that
they must serve to undermine the authority of the Bible; but what were
they to do? The vault of heaven burst, new solar systems appeared; the
earth shrank in a manner highly unbecoming for Christ and Christians.
At last, at the transition from the eighteenth to the nineteenth
century, the Creator was pushed back, the globes were allowed to form
themselves out of a primeval nebula, and in all formations one kept
to matter, forces and laws.

Not everything, to be sure, was settled by this. Though the
theologian could not follow the astronomer on the steep paths among
the stars, he could still console himself that the common people
could not do so either. If it would never become really evident to
the people what the scientific extravagances of the astronomers
signified, then one could simply ignore them, provided one only took
care to keep the almanac clean. The point, then, was to meet science
not in the immeasurable space of the heavens but on the solid earth,
and here to prevent its deviation from the word of revelation in
Scripture. But here too science has taken power from theology.
Geology, the doctrine of the earth's gradual development, is no doubt
still so young that its shape may constantly change; but this much is
already clear, that it conforms not to the biblical doctrine of
creation but to the facts of nature. The oppositions it has to contend
with all fall within science itself; objections from a supernatural
standpoint it does not heed at all. It is, no doubt, only rarely that
theologians seriously engage with geology, but when it occasionally
happens, they tend to betray a colossal misunderstanding. Thus, e.g.,
the theologian Tholuck\footnote{Vermischte Schriften. Zweiter Theil. Hamburg 1839. II. Was ist
das Resultat der Wissenschaft in Bezug auf die Urwelt? p.\ 148
ff.} % *)
speaks strong words, especially on the occasion of a couple of other
theologians' appreciative statements about geology. ``\textit{Die
Erfahrungswissenschaften aller Art, heisst es\ldots} (Bretschneider in
his \textit{Sendschreiben}), \textit{haben fühlbarer und störender in
den Bestand des}
% --- p. 36 ---
% TR: German kept as printed, including "grosze", "entshehenden" (for entstehenden) and "mit den einen" (for dem einen), logged sic in the Danish
\opage{36}\textit{alten theologischen Lehrsystems eingegriffen, als die
spekulative Philosophie\ldots\ Die Geologie will die mosaische
Schöpfung nicht mehr mit den Revolutionen, die unser Erdball erlitten
hat, einstimmig finden. Sie lehrt, unbekümmert, wie der Theologe dabei
zurechtkomme, dass die Erde mehrere grosze Bildungsepochen von
unbestimmbarer aber langer Zeit durchgegangen ist, und dass die ersten
Schöpfungen auf ihr wieder untergegangen sind}''. ``\textit{Die
erhabene Wissenschaft der Astronomie war es\ldots, welche in die
Begriffe des Alterthums von Himmel, Erde, Hölle, Auferstehung,
Gericht, Ende der Welt, die noch zur Zeit der Reformation unverändert
waren, auflösend eingriff --- die Naturforscher und Reisebeschreiber
schilderten die Verschiedenheit der Racen an Gestalt, Farbe und
geistigen Kräften, die durch die Vermischung der Racen entshehenden
Spielarten, und wiesen die grossen und bleibenden Unterschiede unter
ihnen nach, indem sie zeigten, dass diese Differenzen nicht auf
Rechnung des Klimas und der Nahrung, sondern auf \emph{Verschiedenheit
der Grundabstammung} sich gründen müssen. \emph{Blumenbach} sammelte
die Schädel in allen Welttheilen und brachte die Ansicht hiervon in
ein System. In welche Verlegenheit gerieth nun der Theologe? Wenn es
nun nicht mehr \emph{einen} Adam für alle Menschen, sondern einen Adam
für die Kaukasier, einen andern für die Neger, einen dritten für die
Amerikaner, einen vierten für die Malayen, einen fünften für die
Mongolen u.\ s.\ w.\ gegeben hat; wo blieb nun die Dogmatik mit den
\emph{einen} Adam der Bibel, mit der Lehre vom Sündenfall und von der
durch Adam auf alle Menschen gebrachten Schuld, wo nun mit der ganzen
Lehre von der Erbsünde als Folge des Falles und einer von Adam aus
durch Zeugung an alle Menschen gekommener Schwäche? Und ging diese
verloren, wie stand nun die Nothwendigkeit der stellvertretenden
Genugthuung Christi, des zweiten Adams, um die Schuld des ersten Adams
aufzuheben, zu erweisen? Wo blieb nun der Grund der Verdammniss der
Heiden, die nicht von Adam abstammen?}'' ---
% TR: "Tholuch" (sic in the Danish) = Tholuck
Tholuck is extremely indignant at this unreserved statement of
Bretschneider's. ``\textit{So lassen sich}'' --- he exclaims ---
``\textit{die}
% --- p. 37 ---
\opage{37}\textit{Diener der christlichen Kirche im neunzehnten
Jahrhundert nach Entstehung derselben vernehmen. Sieht man, nachdem
eines \emph{Julianus}, \emph{Porphyrius}, \emph{Voltaires} Stimme
verklungen ist, die Pfleger \emph{des christlichen Heiligthums} mit
solchen Reden auftreten, wer muss nicht ausrufen: Herr! Bewahre deine
Kirche vor ihren Freunden und Pflegern, mit ihren Feinden hat es keine
Noth!}''

As Tholuck now reviews the most important geologists from Descartes
down to his own age and goes over paleontology with a fine-tooth comb,
he naturally does not forget to cite a Cuvier, a Leopold von Buch, a
Humboldt and others, but precisely thereby arrives at full assurance of
his own theological result, that the whole thing signifies nothing.
Theology is a true science; but geology, with its many hypotheses, is
not really a science at all. ``\textit{Der christliche Theologe kann in
der That, so lange die geologische Wissenschaft keinen anderen
Standpunkt als den gegenwärtigen einnimmt, in Bezug auf alle jene
Hypothesen nur das stolze Wort gebrauchen, welches jene Königin von
Schweden, Christine, sprach, da sie ihre Krone niederlegte: \emph{non
mi bisogna e non mi basta --- ich bedarf ihrer nicht und sie ist mir
nicht genug}}''. When Tholuck, in order to explain creation, appeals to
the Spirit of God that hovered over the waters, and for the rest
refers to poetry\footnote{Op.\ cit.\ p.\ 172.}, % *)
the only remarkable thing about this is that he cannot himself notice
that he has gone far beyond the limit of knowledge and of science.

As regards the paleontological, the Darwinian hypothesis has in recent
decades so cleared things up in botany and zoology that it may perhaps
be difficult to say of many able specialists on which side they really
stand; only this much is clear, that they can hardly take their stand
on the doctrine of a supernatural creation. E.\ Haeckel, one of the
most eminent adherents of Darwinism, has by learning, acumen and
systematizing talent shown the way to explain the origin of life from
physical causes, conditions and laws, without any interference
whatsoever of a
% --- p. 38 ---
\opage{38}pregiven teleology, and, by demonstrating the necessary
transitional links, has traced the descent of the animals from the
worm to man in one continuous series. His explanation of life we shall
come back to; as for his systematics, zoologists in the future will no
doubt find not only much to acknowledge but also much to correct. A
coherent, progressive development of the organic forms of nature may
very well have its validity without science's therefore succeeding at
once in demonstrating the particular transitional links; to invent
them gives only subjective, not objective, satisfaction. In a treatment
% TR: the print leaves this quotation unclosed (logged sic in the Danish); closed here
of the question ``\emph{Are birds descended from
reptiles}?''\footnote{Skildringer af Dyrelivet i Fortid og Nutid. Kjøbenhavn 1880.} % *)
Dr.\ Chr.\ Lütken has furnished an example of how the thorough,
deliberate researcher conceives his tasks. After following the traces
of the birds as far back as possible in the earth's developmental
history, through the Cretaceous and Jurassic periods down to the
Triassic, he ends with the result that the approximation of the birds
to the reptiles is to be reckoned as little or nothing. ``Perhaps,'' he
says on p.\ 41, ``we should seek the resemblance from the opposite
side, from the \emph{reptiles} over toward the birds? In that case it
is to the past of the reptile class that we must turn, and we shall
not long be in doubt as to what forms will especially claim our
attention; it must be the flying \emph{reptiles}!'' But if we now
consider \emph{the long-snouted flying lizard} (\textit{Pterodactylus
longirostris}), it turns out after all that hope disappoints. ``Only
Cuvier's superior acumen and brilliant eye,'' it says on p.\ 45,
``enabled him, from the imperfect drawings available to him, to work
% TR: the inner quotation opens with » and closes with “ in the print (logged sic in the Danish); set as `...'
out that it was \emph{a flying reptile}, a `bat among the lizards,' but
not a bat in the ordinary sense, not a mammal at all, any more than a
bird.'' Since, on the whole, the flying lizards will not do, the
attempt has been made with the giant lizards (the dinosaurs). ``So
close,'' says our author, following preceding explanations, on p.\ 99,
``had the dinosaurs at long last come to the birds in the
paleontologists'
% --- p. 39 ---
\opage{39}conception that it had to be regarded as exceedingly doubtful
whether one was able to distinguish their footprints from those of
contemporaneous giant ostrich-birds! Then, at about the same time,
some ten years ago, Prof.\ \emph{Cope} in North America and Prof.\
Huxley in England came forward with discoveries and views based on
them, according to which the dinosaurs and the ostriches were supposed
to show such great approximations and agreements, notably in \emph{the
structure of the hind limbs}, that the dinosaurs had to be said,
\emph{at least in this respect}, to stand midway between birds and
ordinary reptiles (crocodiles and present-day lizards).'' It is
notably with a view to \emph{the transformation of the hind limbs}
that the author warns against being finished too early with the
doctrine of descent. ``Perhaps,'' he says on p.\ 111, ``victory has
been prepared in this domain too, and may come before we expect it.
But to send tidings of victory out into the world is not the same as
really having won.'' It is with special regard to the scientific
principle of cognition that we cite this. One does not go back to
miraculous presuppositions because in science one works against
% TR: the print closes this quotation without reopening it after "endog gjærne," (logged sic in the Danish); reopened here
overhasty theories. ``We even readily grant,'' declares our author on
p.\ 5, ``that, were we to make the choice between one of two---either
to assume that all species are created by miracle (be it through a
series of periodically recurring, collective renewals of the whole of
organic nature, or by an infinity of isolated creations), \emph{or} to
assume that the species of the present have developed from those of
the preceding period, and these in turn from those of the period
% TR: comma before capital "Vi" (logged sic in the Danish) rendered as a full stop
before that, etc.---then the choice is not difficult for us. We
naturally prefer the \emph{possible} (more or less probable or
improbable) to what is, scientifically speaking, absurd, even if we do
not see that one has thereby got beyond all difficulties.''

Biology presents a difficulty of a peculiar kind. For not all
physiologists have, as Haeckel has, the courage to deny every
essential difference between the lifeless and the living, a
difference that is supposed to be reducible to certain changes of
molecules and physical forces. It can be expressly demonstrated that
the forces in mechanics, chemical physics and chemistry act
% --- p. 40 ---
\opage{40}by virtue of properties in matter that correspond directly
to the law of inertia and therefore betray no trace of life. That
manifestations of life in the lowest, most imperfect organisms,
monera, e.g., appear to be merely physical effects of physical forces
is an easily explicable deception, which betrays itself when one goes
back from a higher stage of life to lower ones. In Haeckel's
``\textit{Generelle Morphologie der Organismen}'' (p.\ 234, 1st part)
it is said quite naively that the soul of the animal rests only on
representations being awakened in the albumen molecules of the
ganglion cells, i.e., when these molecules themselves begin to
represent; the representation that is awakened by the
\emph{centripetal} fibers is \emph{feeling}, therefore: the molecules
feel; the representation that is awakened in the \emph{centrifugal}
fibers is \emph{will}, therefore: the molecules will; the
representation that is developed under very complicated interactions
of centripetal and centrifugal fibers is \emph{thought}. It is not,
then, the ganglia and brain that merely furnish necessary conditions
for thinking, but the albumen molecules of the ganglion cells are
themselves thinking: the brain is the real subject of thinking. But
however true it is that science can neither demonstrate whence
organic life really comes nor make clear how the punctual interaction
with matter takes place, it is nevertheless evident, and can be
scientifically demonstrated, that a sum of the activity of elementary
forces cannot by any change or transformation whatsoever be converted
into vital activity, that every living organism must rest precisely on
concrete individual unity, unity of difference, original unity of
opposite functions, functions of matter and functions of
selfhood\footnote{Almindelig Videnskabslære. Kjøbenhavn 1880. p.\ 172 ff.}. % *)
Even if, owing to the essential limitation of human knowledge, it is
not possible to reach further, we can still, by reflecting on the
general relation of subjective and objective that runs through the
whole of nature, assure ourselves that the emergence and development
of life rest on peculiar dispositions, which by no means break the
scientific connection by requiring miracles.

% --- p. 41 ---
\opage{41}``\emph{From a higher court to a lower bring not thy cause}''
--- these words of an English physiologist, who will not make
physiology immediately dependent on chemistry, apply especially to the
relation between biology and the elementary sciences. Although physics
and chemistry are such well-founded sciences that no one can doubt
their validity, it is nevertheless undeniable that the foundation on
which they rest, namely atoms and forces, can by no means be said to
be either self-evident or sufficient in and of itself. A system of
world-laws must necessarily be presupposed if a science is to be
possible, and it cannot, as Kant would have it, be our own
consciousness from which the world-laws proceed; but if we think all
consciousness away, then there are no laws, any more than atoms and
forces. Or might the objective perhaps be able to exist, not merely
prior to our finite subjectivity, for that goes without saying, but
prior to all subjectivity, hence before it was objectified? That is
impossible. The laws cannot be without thoughts, since they have the
universal determinations in themselves; but neither can they think
themselves as forces in matter: here a mighty thinking self-unity must
be presupposed, an eternal subjectivity, in which the existence of the
objective world is grounded. Let us but admit that in the original
relation between subjective and objective, between inner and outer,
ideal and real, between the concepts of things and the existence of
things, there is something impenetrable to our limited reason; the
general fundamental relation itself nevertheless shows itself all the
more clearly, the more we think about the reciprocal determination of
things and laws, actuality and possibility, totality and disposition.
If the solar system must presuppose a disposition in which the
objective is originally determined by the subjective, if geological
development must rest on a particular disposition preformed in the
solar system, then the living organisms too must point back to
dispositions that lie in the essence of subjectivity. This much stands
fast, that what is subjective in the organisms cannot be derived from
anything objective whatsoever, but points back to the eternal essence
of subjectivity; but how the infinite subjectivity, which objectifies
laws and things, informs the subjective dispositions into matter
% --- p. 42 ---
\opage{42}is and remains an impenetrable riddle: we stand at an
insurmountable limit.

If, now, it cannot be scientifically demonstrated how life comes forth
in matter and acts; if, as regards its origin, an obscurity must
necessarily remain here; with what right, then, can one demand that
science should be able to give an exact account of what becomes of the
soul when it leaves the body? Strict proofs have been given in the
name of science that the soul is immortal, and in return equally
strict proofs have been given that the psychical itself must dissolve
and vanish at the death of the body. Both attempts are equally
invalid; in both cases science would overstep its proper limit.
Matters stand correspondingly with the scientific treatment of the
question of God's existence. It might indeed seem as though the
subjectivity introduced in the \textit{Videnskabslære}, truth's unity
of knowledge and power, which furnishes the ideal-real presupposition
for all objective actual nature, must be scientifically determinable
as God himself, as the free author of the world, as absolute
personality in original opposition to the impersonal; but that a
characterization of the fundamental relation between subjective and
objective, which runs through all existence, can by no means open to
science any insight into God's supramundane essence, into God as God
with absolute personality, has been expressly emphasized. ``We cannot
speak of spirit and nature as if spirit were outside nature, or nature
outside spirit. The unity at issue here is a self-unity that belongs
essentially to nature, and without which nature cannot possibly be. To
that extent, then, self-knowledge is also knowledge of nature,
self-comprehension comprehension of nature, self-power power of
nature; but, be it noted, not from the standpoint of reality but from
that of ideality. Ideality and reality are no doubt two sides of the
same thing, but these two sides are also two totalities in one
totality. The ideal totality is the spirit of nature; the real
totality is actual nature. Whether one wishes to call this view
theism, naturalism, dualism or monism is quite immaterial. What
matters is exact determinations that can and must be taken up into
the foundation of the
% --- p. 43 ---
\opage{43}science of spirit. The question on which everything depends
is thus this: how far can science reach in determining the fundamental
relation between concept and reality, self-knowledge and knowledge of
another, self-power and power in another, or, expressed otherwise,
between freedom and necessity, spirit and nature''\footnote{Videnskabslære p.\ 227--228.}. % *)

If, now, it is just as impossible for science to prove God's existence
as his non-existence, just as impossible to prove the reality of such
personal attributes as justice, love, mercy in the eternal subjectivity
as to disprove it, just as impossible to prove as to disprove the
supernatural, although a solution of such problems is of decisive
significance for the personal self of man: what then follows from
this? From this it certainly follows that besides scientific
cognition, conditioned by necessity, there must be a religious
cognition, conditioned by freedom. It goes here with science as it
formerly went with religion. When, in the name of religion, the
principle of faith had been made solely valid and ecclesiastical
dogmatics the highest science, the tables turned: reason took power
from faith, and the principle of cognition became scientific. But when
one then, in the name of reason, would make the principle of knowledge
the solely valid, all-embracing principle, the tables turn again, and
we see that there must also be a principle of cognition in faith, that
there must therefore be something that is an object of faith and
cannot possibly become an object of knowledge.

\begin{center}IV.\end{center}

The two principles, that of faith and that of knowledge, are,
according to what we have now seen, not to be regarded as two sides
of the selfsame thing, nor
% --- p. 44 ---
\opage{44}can they be conceived as fragments of a single principle that has
unfortunately cracked and broken to pieces. The two principles each have
their own domain, the principle of faith the supernatural, the principle
of knowledge the natural; and yet both, each in its own way, have the same
domain, namely everything that can be sensed and represented, experienced
and thought. The whole content of faith can, according to what we have
said above, be made an object of knowledge, and conversely the whole
content of knowledge can acquire significance for faith, without faith
thereby being grounded on knowledge or knowledge on faith. It says
(Ps.\ CIV, 14): ``He''---God---``causes grass to spring up for the cattle
and grain for the use of men''; the believer fixes his gaze on the
natural in order to see in it a revelation of the supernatural; the
knower sees in the supernatural a superfluous fancy, since both grass and
grain can grow forth from a natural ground, on natural conditions. The
difference in principle between faith and knowledge cannot, then,
consist in the believer and the knower not understanding each other at
all: but the understanding bends round into contradiction, in that they,
each according to his principle, objectify in different ways: in faith
the personal, and with it the subjective, has decisive preponderance; in
knowledge, by contrast, what matters is the objective, the impersonal.

This shows itself clearly in the conception of laws and miracles.
When, on account of the immediacy with which faith in revelation sets
miracles against natural events in its language, it has been claimed
that such a faith in miracles must necessarily presuppose that the laws
are broken, this is, as far as the principle is concerned, entirely
false; the truth is that faith refers both the laws of nature and the
miracles to God's almighty will. The theological-scientific, i.e.\
unscientific, conception of God's omnipotence has certainly led to
absurd consequences\footnote{Malebranche. De la recherche de la vérité p.\ 390. \textit{Dieu a créé le
monde . . . et produit ainsi tous les effets, que nous voyons arriver,
parce qu’ il a voulu aussi certaines loix, selon lesquelles les mouvemens
se communiquent à la rencontre des corps: et parce que ces loix sont
efficaces, elles agissent et les corps ne peuvent agir.} Om Theologiens
Naturbegreb. Ved Reformationsfesten 1855.}; % *)
but it is, after all, not
% --- p. 45 ---
\opage{45}about theology, it is about religion that we are speaking here. If
science now follows its principle consistently, it must undeniably keep
strictly to unchangeable laws; but since these laws belong at once to
knowledge and to power, and since the unity of power with knowledge may
in truth presuppose will, it is impossible for science to decide whether
the necessity of the laws is not perhaps determined in the last instance
by divine freedom, and so whether the apparent interruption of the laws
of nature by the miracles is not precisely the true expression of the
origin of those laws. Science has, according to its principle and the
objectification determined by it, reason to reject all miracles without
exception, but, be it noted, on the condition that it acknowledges
another principle, insofar as it acknowledges a possibility that it
cannot itself see through. Science is not atheistic insofar as it
declares outright that it does not know the God of religion, for this
God is no scientific quantity; science is not unbelieving because it will
have nothing to do with miracles, for miracles are not objects of
knowledge. When science subjects the accounts of creation and paradise to
a sharp critique, it merely gives an account of the limit of its
cognition, but does not thereby prove that creation and paradise have no
real significance.

If faith, then, has its principle just as well as knowledge, it
follows, of course, that there must be thoughts of faith just as well as
thoughts of knowledge, grounds of faith just as well as grounds of
knowledge, and naturally also a critique of faith just as well as a
critique of knowledge. The critique of faith must then secure the
religious conception of the content of revelation, by not only repelling
the unwarranted interference of knowledge, but also taking care, in
determining the relation between miracle and law of nature, not to come
into contradiction with itself. That faith too must have its critique
shows itself, then, especially in its relation to the miracles. If anyone
will claim that a believer who can accept one single miracle must be
able to accept all, this must no doubt be granted, insofar as faith is
opposed to knowledge; but it by no means follows from this either that
all possible miracle stories must stand on an equal footing in the eyes
of faith, or that the miracles in the
% --- p. 46 ---
\opage{46}New Testament are equally decisive for the life of faith to the same
extent; only the miracle of creation and the miracles of incarnation and
inspiration are fundamental for the critique of faith.

The question of miracles and laws of nature is most intimately
connected with the question of rational reality and spiritual reality,
and this in turn with the question of the objective and its
objectification. Why paradise and the fall cannot possibly assert real
objectivity for themselves in science is obvious; why paradise and the
fall must necessarily assert real objectivity for themselves in faith is
likewise obvious. By assigning to faith and to knowledge each its own
principle of cognition, we seem, then, to make the opposition so strong
that it must necessarily burst the unity of human consciousness; the
believer cannot be a knower, the knower not a believer; and so the whole
ends in an intolerable dualism. ``For the subject standing outside,''
% TR: the print opens „ before "Om Holbergs" in this note and never closes it; closed here after the title
says Strauss\footnote{Cf.\ ``Om Holbergs Kirkehistorie og Theologi'' by the present author.
Copenhagen 1867. p.\ 71 ff.}, % *)
``revelation must prove by external objective signs that it is
revelation; it must document its origin from the absolute intelligence
by showing, in the domain that the finite spirit can check, how it
surpasses the latter, which happens through \emph{prophecies}; it must
prove its origin from the absolute power by dominion over the power of
nature, through \emph{wonders}. In order to be carried down to posterity,
revelation must be preserved in \emph{holy scriptures}, but these must,
so that one can be sure that revelation is kept undistorted in them, be
inspired by the author of revelation. Yet this whole apparatus for making
revelation secure and accessible to us still lies throughout on the
objective side; there is still lacking the most important thing, the
subject's appropriation of the revealed content. The subject appropriates
revelation through the \emph{interpretation of Scripture}, but how is it
to assure itself that the alleged revelation is really divine? Through
the wonders and prophecies that accompanied it when it took place. Good,
but how will one
% --- p. 47 ---
\opage{47}be able to know that these are real wonders and prophecies? From the
testimony of Scripture. But how does one know that this is true? From
its being inspired by the infallible God. But that Scripture is inspired
by God, how does one know that? From the \emph{inner testimony of the
Holy Spirit}, which, when we read Scripture, recognizes in it its own
work. Good, but what is now to convince us that this testimony in us
really comes from the Holy Spirit, and not from our own spirit, or even
from a seducing spirit outside us?---Here the thread of the orthodox
system snaps... He who does not have the spirit within him has it
outside him; he who cannot find the determining principle in himself
seeks it in a revelation; he who is not yet ripe for the faith of
reason remains with faith in revelation.'' Relying on what he calls the
faith of reason, Strauss points to a chasm between two classes of human
society, the knowers and the people, the non-philosophizing in the
higher as well as in the lower ranks, a chasm that will perhaps never be
filled. ``So let,'' he says, ``the believer then let the knower, just as
the latter the former, go his way in peace; we let them keep their
faith, so they must also let us keep our philosophy; and if the overly
pious should succeed in excluding us from their church, we would count
it a gain; enough false attempts at union have now been made; only a
sharpening of the oppositions can lead further.''

That the Straussian separation between believers and knowers is
thoroughly false can be seen from what we have said above. Reason, if we
keep to the principle, has neither any religion nor a substitute for
religion to offer us; every religious object, even an Egyptian Apis,
presupposes a believing subjectivity and, in all its naturalness, points
to the supernatural. If this supernatural has nothing higher, nothing
subjective or personal, in it, it must consistently consume man's
personal self, and faith then shows itself active from the negative
side. Thoroughgoing nihilism, which dissolves all faith, dissolves all
knowledge too. This is seen clearly from ``The Philosophy of the
Unconscious''\footnote{Philosophie des Unbewussten. Versuch einer Weltanschauung von
Ev.\ Hartmann, Dr.\ phil.\ Berlin 1869.}. % *)
Right from Schelling's first
% --- p. 48 ---
\opage{48}youthful appearance with the intellectual intuition of the absolute
down to the most recent time, speculative philosophy, in ambiguous
understanding with natural science, has labored to make the absolute
being clear to itself. It has, however, been difficult for the
intelligent, co-philosophizing public to keep up and to find out how
matters really stood with this absolute. Now it was presented as a
logically omniscient deity: God in the Logic before he created the
world; now again as an unopened, obscure principle, the indwelling
ground and cause of all existence. Now at last the ambiguity has been
removed; the countenance of the god of truth is seen completely
unveiled. The riddle of the world is that an unconscious willing has from
the beginning driven a likewise unconscious representation from
% TR: the quotation opened at "in itself" has no closing mark in the print (4 „ against 3 “ on p. 48); closed here after the second "gar Nichts"; "gar Nichts" set as inner quotation
non-being into being, for, it says on p.\ 607, ``in itself representation
had no drive and no interest to pass from non-being over into being;
consequently, before the entry of will no representing was actual;
consequently, before the coming-to-be of the world there was neither
willing nor representing, i.e.\ `\textit{gar Nichts}'. As long as willing
lasts, so long will the process and its appearance in consciousness, the
world, last; when, therefore, there is one day to be no world any more,
then there may no longer be any willing or any representation either,
since unconscious representation only ever becomes actual insofar as the
interest of will demands it, i.e.\ there will again be `\textit{gar
Nichts}'.'' With this, then, the following description of the deity also
agrees (p.\ 606): ``The All-One Unconscious is omniscient and all-wise,
and consequently cannot become wiser than it already is; it has, as
Aristotle too says, no memory, so it can learn nothing from the
experiences it otherwise had in the world. Consequently, when the world
has once ceased to be, it has remained exactly the same as it was before
the creation of the world; as blessed as it formerly was, it is now
again, neither more nor less; the world process can in no way help it to
greater blessedness than it formerly had.'' If the deity is blessed in
nothing, it follows of itself that the soul too must become blessed in
nothing. The hope of an individual continuance of the soul after death
is declared outright to be an illusion.
% --- p. 49 ---
\opage{49}``With that the main nerve of the Christian promises is cut, for at
bottom man cares only about his own dear I; what good to me is the
greatest future blessedness if I cannot perceive and enjoy
it!''\footnote{For Ide og Virkelighed. Vol.\ 1875. Hvorledes man danner sig en
Verdensanskuelse. p.\ 12--14.} % *)

In nihilism not only faith but knowledge too perishes. That out of
nothing itself, pure nothing, a representing and willing should have
formed themselves, and by their union a real world, which will in turn
dissolve itself into the infinite nothing, is surely as meaningless as
possible. A real existence, an original unity-in-opposition of subjective
and objective, and, as far as man is concerned, an essential difference
between impersonal objectifying in knowledge and personal objectification
in faith, must necessarily be presupposed. We shall try to illustrate
this by some examples. That the biblical account of the creation of the
world does not agree with the scientific conception of the world of
things contains, when we look more closely, no contradiction. When one
and the same object is objectified by two persons under equal conditions
with the same sense, e.g.\ sight, and yet presents itself differently,
% TR: the print reads "om blaat" (for "som blaat"); rendered "as blue"
to the one as blue, to the other as gray, there must naturally be a fault
in the organ of sight either in the one or in the other, if not in both;
but when the one can judge only by sight, the other by contrast by his
own feeling, one cannot demand full agreement, even if the object in and
for itself is assumed to be the same. Science, which dwells on natural
elements, atoms, primal nebula, physical causes, conditions and laws,
does not reach spirit; its impersonal objectification, however rational
it may be, must stop at the externally real. In religion, by contrast,
objectification is originally determined by the peculiar co-operation of
spirit. The earth, the sun, the moon are indeed conceived as they
immediately show themselves to sense and understanding, but what is
decisive lies in the assurance that the spirit of faith awakens about
their origin and their relation to
% --- p. 50 ---
\opage{50}the subjective, the personal. It must be expressly emphasized that the
scientific presentation of the universe cannot possibly come into
conflict with the religious one, unless science goes beyond its limit,
abolishes the imagination of revelation, and itself undertakes to decide
how matters stand in principle with creator and creature. Nor can the
religious conception of a creation in six days come into conflict with
the doctrine of a successive development, unless religion will give up
its limitation in principle, set itself in the high seat of science, and
show clearly how millions of years shrink to six days.

``But,'' one objects, ``if a settlement between religion and science is
to be grounded on so principled a difference in the mode of
objectification that science gives real objectivity, freed from all
extraneous interferences, as it is, while religion gets its peculiar world
through an objectification that transforms reality by mixing in the
personal, then the whole thing is confused. What, pray, becomes of this
personal when we speak plainly? The mythical! If expressions such as
personal objectivity, spiritual reality, reality of revelation are
nothing other than paraphrases of what we already know under the name of
the mythical, then nothing at all is gained hereby toward a real
illumination of the self-validity of religion over against science.'' We
shall test this objection by casting a glance at paradise. That
paradise, in every form, must in science be an illusory reality, a
delusion of the imagination, a myth, cannot, according to what we have
said above, be contradicted; but we have, after all, granted faith its
principle, so that it should not stand under the tutelage of knowledge,
but have the right to assert its own view of the world and to apply its
own critique. Let us then, with the principle of faith before our eyes,
consider paradise from the standpoint of revelation, of myth and of
poetry. Revelation begins with a setting in which the deity, in a
sensuous-supersensuous presence, lives together with man; this paradisal
reality is spiritual reality. The immediate fusion of natural and
supernatural, however, already belongs to the past when paradise becomes
an object of
% --- p. 51 ---
\opage{51}faith. Since the past cannot be seen immediately but must be viewed in
the reflecting mirror of the imagination, it becomes uncertain whether
faith can find the given signs, from the divine as well as from the
human side, so satisfying that it must hold fast in the paradisal
phenomena to essential features of a true revelation, or whether, as
reflection is set in motion, an uncertainty sets in about what is divine
and what is human, how much is supernatural, how much is natural. In the
latter case revelation is transformed into myth. The mythical has many
changing shapes, and furnishes a foundation for poetry and art. The
paradise with which Ovid begins his Metamorphoses is evidently merely
poetic; the paradise of the Persian religion is mythical; but why is the
paradise of the Hebrews and the Christians neither merely poetic nor
mythical, but a real revelation? That must be decided in principle in
faith, not in knowledge. According to the principle of knowledge all
paradises are mythical, the Mosaic-Christian just as well as the Persian
and the Muhammadan. The paradisal reality of faith is spiritual reality;
the faith that in the deepest sense recognizes the personal relation
between God and man must therefore, in the name of revelation, recognize
the true, spiritually real paradise.

If we are to examine in particular the activity of Christ and of the
devil, the same principle applies, the principle of faith. That there
really lived in Palestine a man by the name of Jesus who said that he was
the Christ, the only-begotten Son of God, science will not deny; Jesus of
Nazareth is really a historical personality, and as such has a claim to
historical faith. But historical faith (\textit{fides historica}) is
different in principle from religious faith (\textit{fides religiosa});
Jesus as God's Son become man, Jesus as the virgin-born, as the
wonder-worker, as the risen and ascended one, in short: Jesus as the
Christ is the object not of merely historical but of religious faith.
What holds of Christ holds also of the devil. The devil himself, it is
true, is not yet incarnate---the Antichrist, after all, has not yet
come---yet it is seen from what we have said above that the scientific
critique that begins by
% --- p. 52 ---
\opage{52}dissolving the personality of the devil must consistently end by
dissolving the personality of Christ; from which, then, the converse also
follows, that the faith that in principle will hold fast to the personal
Christ must also acknowledge the personal devil.

If one will now conclude from this that, since faith and knowledge are
so diametrically opposed, the believer cannot possibly be a knower, and
the knower not a believer, then the conclusion is utterly false. If the
objects of faith are so different in principle from the objects of
knowledge that one can be a believer without being a knower, because the
religious objects are personal, the scientific ones by contrast
impersonal, then it also follows from this that the believer can be a
knower, and the knower a believer, without coming into contradiction with
himself. An astronomer who would at once embrace both the Ptolemaic and
the Copernican world system certainly comes into contradiction with
himself, for both systems are results of the same kind of
objectification: if the one is true, the other must necessarily be
untrue. He, on the other hand, who in his capacity as astronomer
embraces the Copernican system, but in his capacity as believer the
Mosaic account, does not come into contradiction with himself; for in
aiming through the natural at the supernatural and objectifying
personally in imagination, he cannot possibly objectify scientifically.
When one person holds fast in faith that man was created on the sixth
day, while another believes that he was created only six million years
after the first formation of the earth, they are in real contradiction
with each other, for they move within the same sphere; but when a
believer asserts that man is created in the image of God, while a
zoologist asserts that man has developed from the same basic form as the
ape, then there is no contradiction in this; for they are speaking of
utterly different things. ``That God Jehovah formed man, dust of the
earth, and breathed into his nostrils the breath of life'' (Gen.\ II, 7)
concerns not the natural development but solely the origin from God,
which can be made intuitable only in personal objectification; man's
development from a natural basic form decides nothing about creation.
If, on the other hand, the believer will not be content to emphasize the
% --- p. 53 ---
\opage{53}divine act of creation, of which the inbreathing of the breath of life
is a significant symbol, but ventures to determine the course of
development, he violates his principle and is thereby already in the
wrong. If, conversely, the scientific researcher, who according to the
principle of knowledge objectifies to himself the stepwise development of
nature, will deny the moment of creation lying in the original
disposition, which he does not know at all, he oversteps his limit and
violates his principle. Here there is on both sides a contentious
insistence on being right which harms religion and science in equal
measure.

Should one perhaps object that it is only in certain special cases that
one can appeal to a mode of objectification different in principle, and
to a correspondingly different kind of reality, in order to remove the
conflict between faith and knowledge, but that this way out is far from
sufficing in the most difficult, most important cases, then let us once
more cast a glance at Christ's resurrection and ascension. If it is
really the man Jesus, and not a general idea of Christ or a merely
spiritual Christ, who has risen from the dead and ascended into heaven,
then the object of faith cannot possibly be detached from the object of
knowledge without the whole reality having to become mythical. He who
rises from the grave and shows himself to the disciples is, after all,
the crucified one, he whom all, unbelievers and believers alike, could
see with sensuous eyes and grasp with sensuous hands, and so objectify to
themselves both in the manner of knowledge and in that of faith; but
after the resurrection his form is as if wholly transformed, real and yet
super-real; Mary Magdalene cannot recognize him (John XX.\ 15--17), nor
% TR: the print has "(Luk. XXIV;" with no closing parenthesis; closed here
can the two who went to Emmaus (Luke XXIV); he goes through closed doors,
and the disciples think they see a spirit; but Thomas, who wants tangible
certainty, gets it and believes (John XX; 19--31). Of a scientific
presentation of such events there can be no question here at all. The
objectivity is, like religion itself, the revealing of what is
incomprehensible in personality, spiritual reality in contemplable form,
an object of faith, not of knowledge.

That faith and knowledge, each with its principle, do not stand in
contradiction with each other must surely be clear from the foregoing;
% --- p. 54 ---
\opage{54}but faith and knowledge are, after all, heterogeneous cognitions: the
question then becomes how they can be united in the unity of human
consciousness. To this the following must be answered: science, having
become fully conscious of its principle, its right, its freedom, cannot
see in religion as divine revelation any enemy that will intervene
disruptively in its investigations; even if scientific cognition cannot
fathom personality itself, there must nonetheless be reason in all
personal thinking. For its part, religion cannot then stand hostile
toward science either; even if, for personal reasons, it refers the laws
of existence to God's will and lets the natural be subordinate to the
supernatural, this subordination cannot possibly mean any breach of the
laws. When it is said that Christ was under the law in order to fulfill
the law, this holds not of the Mosaic law alone but of the law of the
% TR: the print reads "Skiøndt" (for "Skjøndt"); no effect on the sense
world, the law of nature. Although religion, with its principle of
personality, is not dependent on science, it is not therefore excluded
from science. The principle of religion is also the principle of
morality; but the principle of morality works in society on such natural
conditions that in its universal validity it must satisfy the demands of
reason, and for that reason it has been confused with the principle of
reason itself.

Religion too, or rather the doctrine of religion, has need of
scientific illumination, only this illumination must not give itself the
appearance of scientific grounding. ``The philosophy of religion does not
stop at general determinations of the limits of faith and knowledge; it
% TR: the print reads "ude" (for "uden"); rendered "without"
cannot possibly let religion develop its content without letting a whole
system of concepts of faith come to development. But the religious
content of faith is, according to the principle, a suprascientific
content; here, then, the contradiction arises that a suprascientific
content is taken up into science. If the form of science were now really
confused with the true, original form of the content and introduced and
held fast in place of it, the contradiction would certainly be
insoluble. But this is by no means the case. Scientific reflection has
precisely the
% --- p. 55 ---
\opage{55}flexibility that it brings the form of knowledge to reflect another and
opposite form. For just as the clear day can bring the eye to forget the
gleam of light and look at the objects, so too reflection can bring
thought, in knowledge, to look at faith and forget knowledge. The
philosophy of religion is so far from forcing an alien form on the
religious content that its development of form consists, on the
contrary, in freeing that content from admixtures of immediately
% TR: the print reads "o, s. v." (comma after "o"); rendered "etc."
individual representations, images, feelings, expressions of will, etc.,
with which it is always grown together in the life of faith itself.
Insofar as reflection, as regards the form as well as the content,
cancels knowledge through knowledge, in that it aims to clarify not a
scientific but a suprascientific truth, the relation between truth and
science is essentially reversed, and the philosophy of religion is in
this respect to be characterized as an inverted
science.''\footnote{Religionsphil.\ p.\ 8--11.} % *)

\emph{You shall know the truth}, and \emph{the truth shall make you
free}.
% end of Nielsen's text: a short double rule is printed below the footnote on p. 55; no signature or date

\end{document}
